% Copyright 2026 The Android Open Source Project
%
% Licensed under the Apache License, Version 2.0 (the "License");
% you may not use this file except in compliance with the License.
% You may obtain a copy of the License at
%
%      http://www.apache.org/licenses/LICENSE-2.0
%
% Unless required by applicable law or agreed to in writing, software
% distributed under the License is distributed on an "AS IS" BASIS,
% WITHOUT WARRANTIES OR CONDITIONS OF ANY KIND, either express or implied.
% See the License for the specific language governing permissions and
% limitations under the License.
%
% Build with:  tectonic doc/spec.tex

\documentclass[11pt,a4paper]{article}

\usepackage[T1]{fontenc}
\usepackage[utf8]{inputenc}
\usepackage{lmodern}
\usepackage{microtype}
\usepackage[margin=2.6cm]{geometry}
\usepackage{amsmath}
\usepackage{amssymb}
\usepackage{amsthm}
\usepackage{booktabs}
\usepackage{longtable}
\usepackage{array}
\usepackage{enumitem}
\usepackage{xcolor}
% Long Java identifiers appear inline constantly; let them hyphenate rather
% than punch through the right margin.
\usepackage[htt]{hyphenat}
\usepackage[colorlinks=true,linkcolor=black,urlcolor=blue,citecolor=blue]{hyperref}

%%% ------------------------------------------------------------------ macros

\newcommand{\Sc}{\mathbb{S}}
\newcommand{\Om}{\Omega}
\newcommand{\bo}{\bot}

\newcommand{\Exact}{\mathsf{Exact}}
\newcommand{\Fill}{\mathsf{Fill}}
\newcommand{\Wrap}{\mathsf{Wrap}}
\newcommand{\Weight}{\mathsf{Weight}}

\newcommand{\pad}{\mathsf{pad}}
\newcommand{\wdt}{\mathsf{wdt}}
\newcommand{\hgt}{\mathsf{hgt}}
\newcommand{\pri}{\mathsf{pri}}

\newcommand{\Horz}{\mathsf{H}}
\newcommand{\Vrt}{\mathsf{V}}

\newcommand{\leaf}{\mathsf{leaf}}
\newcommand{\ttext}{\mathsf{text}}
\newcommand{\bbox}{\mathsf{box}}
\newcommand{\rrow}{\mathsf{row}}
\newcommand{\ccol}{\mathsf{col}}
\newcommand{\crow}{\mathsf{crow}}
\newcommand{\ccolp}{\mathsf{ccol}}
\newcommand{\fflow}{\mathsf{flow}}

\newcommand{\miw}{\mathsf{miw}}
\newcommand{\mih}{\mathsf{mih}}
\newcommand{\gone}{\mathsf{gone}}
\newcommand{\vis}{\mathsf{vis}}

\newcommand{\Start}{\mathsf{Start}}
\newcommand{\Center}{\mathsf{Center}}
\newcommand{\End}{\mathsf{End}}
\newcommand{\Top}{\mathsf{Top}}
\newcommand{\Bottom}{\mathsf{Bottom}}
\newcommand{\SBetween}{\mathsf{SpaceBetween}}
\newcommand{\SEvenly}{\mathsf{SpaceEvenly}}
\newcommand{\SAround}{\mathsf{SpaceAround}}

\newcommand{\Mm}{\mathcal{M}}          % measure judgement
\newcommand{\Rr}{\mathcal{R}}          % axis resolution
\newcommand{\Aa}{\mathcal{A}}          % arrangement
\newcommand{\Pp}{\mathcal{P}}          % padding accumulation
\newcommand{\Dw}{\Delta_{\mathrm{w}}}
\newcommand{\Dh}{\Delta_{\mathrm{h}}}
\newcommand{\Iw}{\iota_{\mathrm{w}}}
\newcommand{\Ih}{\iota_{\mathrm{h}}}
\newcommand{\Dimw}{\mathsf{dim}_{\mathrm{w}}}
\newcommand{\Dimh}{\mathsf{dim}_{\mathrm{h}}}
\newcommand{\base}{\mathsf{base}}
\newcommand{\val}{\mathsf{val}}
\newcommand{\wt}{\mathsf{wt}}

\newcommand{\clampL}{\mathord{\uparrow}}
\newcommand{\clampH}{\mathord{\downarrow}}

\newcommand{\lo}[1]{\texttt{#1}}       % a Lean identifier
\newcommand{\jv}[1]{\texttt{#1}}       % a Java identifier
\newcommand{\file}[1]{\texttt{#1}}

\DeclareMathOperator{\clamp}{clamp}
\DeclareMathOperator{\len}{len}

\theoremstyle{definition}
\newtheorem{definition}{Definition}[section]
\theoremstyle{plain}
\newtheorem{property}{Property}[section]
\theoremstyle{remark}
\newtheorem*{remark}{Remark}
\newtheorem{warning}{Caveat}[section]

\setlist[description]{style=nextline,leftmargin=1.2em}

%%% ------------------------------------------------------------------ title

\title{\bfseries A Formal Specification of the\\ RemoteCompose Layout Algorithm\\[4pt]
  \large Core Managers, Adaptive Containers, Extended Subsystems, and Optimizations}
\author{Machine-checked in Lean~4 \\
  \texttt{compose/remote/specification/conformance/lean}}
\date{}

\begin{document}
\maketitle

\begin{abstract}
\noindent
This document is a mathematical specification of the complete layout engine
implemented by the RemoteCompose player. It formalizes nine layout managers:
the core containers (\jv{Box}, \jv{Row}, \jv{Column}),
line-breaking and adaptive containers (\jv{Flow}, collapsible \jv{Row} and \jv{Column}),
and specialized containers (\jv{FitBox}, \jv{State}, \jv{Image}).
Beyond container sizing and placement, it specifies the modifier pipeline
--- including visibility modes (\jv{VISIBLE}, \jv{INVISIBLE}, \jv{GONE}),
typographic baseline alignment, visual translation offsets, scrollable viewports,
affine graphics layers, and $z$-ordering --- as well as hierarchical window coordinates,
interactive pointer hit-testing, temporal layout animations, dynamic
expression evaluation, and document-level layout optimization passes.
It is the prose counterpart of a Lean~4 development that defines the entire engine
as total recursive functions, proves core structural properties, and validates
computed rectangles and component visibility against the project's conformance
gold corpus. Every rule is traceable to the reference Java implementation and its
corresponding Lean definition.
\end{abstract}

\tableofcontents

%%% ==================================================================
\section{Introduction}

\subsection{Purpose}

RemoteCompose documents carry a component tree that a player must lay out
identically on every platform. The Java player in \file{remote-core} is the de
facto reference, but its behaviour is defined only by its own source code:
several rules are order-sensitive, two of the managers that look like mirror
images are not, and a handful of arithmetic paths overflow their container by
construction. Reimplementing the player --- in TypeScript, C++, or anywhere else
--- therefore means reverse-engineering those rules from the Java, which is
error-prone and unauditable.

This specification exists to make the algorithm an object one can read, cite and
check. It has three layers:

\begin{enumerate}[leftmargin=1.6em]
  \item \textbf{This document}, which states the algorithm as a system of
    mathematical definitions.
  \item \textbf{A Lean~4 development}, which states exactly the same algorithm
    as a total function \lo{measure} and proves properties of it. The Lean text
    is the authority: this document is generated from it by hand, and any
    disagreement is a bug in this document.
  \item \textbf{A conformance harness}, which replays the project's existing
    gold corpus against the Lean function and checks the resulting rectangles
    inside the Lean kernel.
\end{enumerate}

\subsection{Relation to the reference implementation}

The Java sources referred to throughout live under
\begin{center}
\file{compose/remote/remote-core/src/main/java/} \\
\file{androidx/compose/remote/core/operations/layout/}
\end{center}
and the relevant files are
\begin{itemize}[leftmargin=1.6em,itemsep=0pt,topsep=4pt]
  \item \jv{LayoutManager.java}, \jv{LayoutComponent.java} and
    \jv{Component.java};
  \item \jv{BoxLayout.java}, \jv{RowLayout.java}, \jv{ColumnLayout.java},
    \jv{CollapsibleRowLayout.java}, \jv{CollapsibleColumnLayout.java},
    \jv{FlowLayout.java} and \jv{CollapsiblePriority.java}, in \jv{managers/};
  \item \jv{BaseModernMeasurePolicy.java}, in \jv{managers/policies/}.
\end{itemize}
Two files outside that tree also matter, because they decide what a document
\emph{means} before layout ever runs: \jv{DefaultComponentParsers.java} fixes the
per-kind alignment defaults of Caveat~\ref{sec:defaults}, and
\jv{DefaultModifierParsers.java} resolves a bare \jv{weight} or
\jv{collapsiblePriority} against the axis of the enclosing container. Both live
under
\begin{center}
\file{compose/remote/remote-creation-core/\ldots/json/}.
\end{center}

Measurement is dispatched through a \emph{policy chain}. \jv{LayoutManager}
holds a static array
\[
  \jv{POLICIES} = [\,
    \jv{Legacy},\;
    \jv{BaseModern},\;
    \jv{InsetWrap},\;
    \jv{InlineExpression},\;
    \jv{EnforceConstraints}\,]
\]
indexed by the document's measure version, and each entry from index~1 onwards
extends its predecessor by flipping one boolean. Since
\jv{LayoutManager.DEFAULT\_MEASURE\_TYPE} is \jv{ENFORCE\_CONSTRAINTS}, the
behaviour that matters is \jv{BaseModernMeasurePolicy.measure} evaluated with
\[
  \begin{gathered}
    \jv{shouldApplyInsetWrap} = \jv{shouldUpdateComponentValues} \\
    = \jv{shouldEnforceConstraints} = \textsf{true},
  \end{gathered}
\]
and that is what this document specifies. Line numbers cited as
``policy~$\ell_1$--$\ell_2$'' refer to \jv{BaseModernMeasurePolicy.java}.

\subsection{Scope}

\begin{table}[h]
\centering
\small
\begin{tabular}{@{}ll@{}}
\toprule
\textbf{Subsystem} & \textbf{Formally Specified and Verified Components} \\
\midrule
\textbf{Core Containers} & \jv{BoxLayout}, \jv{RowLayout}, \jv{ColumnLayout} \\
\textbf{Adaptive Containers} & \jv{CollapsibleRowLayout}, \jv{CollapsibleColumnLayout}, \jv{FlowLayout} \\
\textbf{Extended Containers} & \jv{FitBoxLayout}, \jv{StateLayout}, \jv{ImageLayout} (aspect ratio scaling) \\
\textbf{Sizing Modifiers} & \jv{width}, \jv{height}, \jv{widthIn}, \jv{heightIn}, \jv{padding}, \jv{weight} \\
\textbf{Arrangement} & \jv{spacedBy}, all alignments (\jv{Start}, \jv{Center}, \jv{End}, \jv{SpaceBetween}, \ldots) \\
\textbf{Extended Modifiers} & \jv{visibility} (\jv{VISIBLE}, \jv{INVISIBLE}, \jv{GONE}), \jv{alignByBaseline}, \\
 & \jv{offset}, \jv{scroll} viewports, \jv{graphicsLayer}, \jv{zIndex} \\
\textbf{Window Hierarchy} & Hierarchical coordinate accumulation (\jv{getLocationInWindow}) \\
\textbf{Interaction} & Geometric pointer inclusion, hit-testing soundness (\jv{hitTest}) \\
\textbf{Layout Animations} & Temporal progress $\alpha(t)$, bounds interpolation, start/end continuity \\
\textbf{Expressions} & Dynamic expression AST, runtime variable binding, clamping bounds \\
\textbf{Optimizations} & Relayout boundaries, memoized cache, flat-pass generation \\
\textbf{Text \& Autosizing} & \jv{CoreText}, proportional font scaling, autosize bounds \\
\textbf{Repaint Dynamics} & Invalidation cycles, Choreographer attractor, fix verification \\
\bottomrule
\end{tabular}
\caption{Scope of this specification.}
\end{table}

Terminal components include both generic leaves (spacers, canvases, images)
abstracted as intrinsic content sizes and \jv{CoreText} components with nominal
or dynamic autosizing metrics.

%%% ==================================================================
\section{Preliminaries}

\subsection{Scalars}

\begin{definition}[Scalar domain]
Let $\Sc = \mathbb{Q}$.
\end{definition}

The reference implementation computes in IEEE-754 single precision. Modelling
that faithfully would make every statement below carry a rounding hypothesis and
would make the conformance checks undecidable in practice. We therefore work in
$\mathbb{Q}$, which is exact, decidable and closed under the four operations the
algorithm uses. The gap between the two is absorbed by comparing against the
gold corpus with the same tolerance the existing harness uses,
$\varepsilon = \tfrac12$ device pixel (Section~\ref{sec:conformance}).

\begin{definition}[Float sentinel]
\[
  \Om \;=\; 2^{128} - 2^{104} \;=\; 340282346638528859811704183484516925440 .
\]
\end{definition}

$\Om$ is the exact value of \jv{Float.MAX\_VALUE}. It is not a symbolic
infinity: \jv{LayoutComponent.computeModifierDefinedWidth} returns it as an
ordinary finite number for a \jv{FILL} axis, and the caller then takes a
\jv{Math.min} against the incoming maximum. Arithmetic on $\Om$ is therefore
observable and must be exact, which is why it is written out in full.

\subsection{Geometry}

A \emph{size} is a pair $\sigma = (\sigma_w, \sigma_h) \in \Sc^2$; a
\emph{rectangle} is a quadruple $(x, y, w, h) \in \Sc^4$.

\begin{definition}[Padding]
A padding is $p = (p_\ell, p_r, p_t, p_b) \in \Sc^4$, with the abbreviations
\[
  p_{\mathrm{h}} = p_\ell + p_r, \qquad p_{\mathrm{v}} = p_t + p_b .
\]
\end{definition}

\begin{definition}[Constraints]
A constraint box is $\kappa = (w^-, w^+, h^-, h^+) \in \Sc^4$, read as
``the measured width must lie in $[w^-, w^+]$ and the measured height in
$[h^-, h^+]$''. We write
\[
  \kappa_{\mathrm{tight}}(w,h) = (w,w,h,h), \qquad
  \kappa_{\mathrm{loose}}(w,h) = (0,w,0,h).
\]
\end{definition}

\begin{definition}[One-sided clamps]
For $x \in \Sc$ and $o \in \Sc \cup \{\bo\}$,
\[
  x \clampL o = \begin{cases} x & o = \bo \\ \max(x, o) & o \in \Sc \end{cases}
  \qquad\qquad
  x \clampH o = \begin{cases} x & o = \bo \\ \min(x, o) & o \in \Sc. \end{cases}
\]
\end{definition}

$\bo$ models the two distinct ``absent'' encodings the wire format uses:
\jv{Float.NaN} for a missing \jv{FILL} fraction, and $-1$ for a missing
\jv{widthIn} bound. Both are read back as $\bo$ here.

%%% ==================================================================
\section{Syntax}

\subsection{Dimensions}

\begin{definition}[Dimension descriptor]
\[
  d \;::=\; \Exact(v) \;\mid\; \Fill(\varphi) \;\mid\; \Wrap
        \;\mid\; \Weight(\omega),
  \qquad v, \omega \in \Sc,\; \varphi \in \Sc \cup \{\bo\}.
\]
\end{definition}

These are the four inhabited cases of
\jv{DimensionModifierOperation.Type}, whose full wire encoding is
\[
  \begin{array}{llll}
    0\ \jv{EXACT} & 1\ \jv{FILL} & 2\ \jv{WRAP} & 3\ \jv{WEIGHT} \\
    4\ \jv{INTRINSIC\_MIN} & 5\ \jv{INTRINSIC\_MAX} & 6\ \jv{EXACT\_DP} & \\
    7\ \jv{FILL\_PARENT\_MAX\_WIDTH} & 8\ \jv{FILL\_PARENT\_MAX\_HEIGHT}. & &
  \end{array}
\]
Cases 4--8 are out of scope; \jv{EXACT\_DP} is folded into $\Exact$ by the
density conversion that happens before layout.

Two projections are used throughout:
\[
  \val(d) = \begin{cases}
    v & d = \Exact(v) \\
    \varphi & d = \Fill(\varphi) \\
    \bo & d = \Wrap \\
    \omega & d = \Weight(\omega)
  \end{cases}
  \qquad
  \wt(d) = \begin{cases}
    \omega & d = \Weight(\omega) \\
    0 & \text{otherwise.}
  \end{cases}
\]

\begin{definition}[Intrinsic base of a dimension]
For $b \in \{\min, \max\}$,
\[
  \base_b(\Exact(v)) = v, \quad
  \base_b(\Fill(\_)) = \begin{cases} 0 & b = \min \\ \Om & b = \max \end{cases}, \quad
  \base_b(\Wrap) = \base_b(\Weight(\_)) = 0 .
\]
\end{definition}

\subsection{Modifiers}

\begin{definition}[Modifier chain]
\[
  m \;::=\; \pad(p) \;\mid\; \wdt(d, \iota) \;\mid\; \hgt(d, \iota)
        \;\mid\; \pri(o, v),
  \qquad \iota \in (\Sc\cup\{\bo\})^2 \cup \{\bo\},
\]
\[
  o \in \{\Horz, \Vrt\}, \qquad v \in \Sc,
\]
and a chain is a finite \emph{ordered list} $\mu = m_1 \cdot m_2 \cdots m_k$,
written $\varepsilon$ when empty.
\end{definition}

The $\iota$ component of a size modifier carries the associated
\jv{widthIn} / \jv{heightIn} bounds as a pair $(\iota^-, \iota^+)$, each
possibly $\bo$.

A $\pri(o, v)$ is a \jv{CollapsiblePriorityModifier}. It affects nothing unless
the node's parent is a collapsible manager, and it names the axis it applies to
--- a priority declared for one axis is invisible to the other. Its reading is
$P_o$, defined in Section~\ref{sec:crow}.

Purely decorative operations --- \jv{BackgroundModifierOperation},
\jv{BorderModifierOperation}, \jv{ZIndexModifierOperation}, clipping, click
handling --- take no part in measurement and are simply absent from $\mu$.

\begin{warning}
The chain is a \emph{list}, not a record, and the order is observable. An
earlier draft of this specification modelled modifiers as a flat record of
optional fields; it disagreed with the gold corpus on eight documents. See
Section~\ref{sec:orderquirk}.
\end{warning}

\subsection{Nodes}

\begin{definition}[Component tree]
\[
  \begin{aligned}
    n \;::=\;\;
      & \leaf(\mu, \sigma)
      && \text{a spacer or fixed atom} \\
    \mid\;\; & \ttext(\mu, \theta)
      && \text{\jv{CoreText} (nominal or autosizing)} \\
    \mid\;\; & \bbox(\mu, a_{\mathrm{h}}, a_{\mathrm{v}}, \vec{n})
      && \text{\jv{BoxLayout}} \\
    \mid\;\; & \rrow(\mu, A, a_{\mathrm{v}}, s, \vec{n})
      && \text{\jv{RowLayout}} \\
    \mid\;\; & \ccol(\mu, a_{\mathrm{h}}, A, s, \vec{n})
      && \text{\jv{ColumnLayout}} \\
    \mid\;\; & \crow(\mu, A, a_{\mathrm{v}}, s, \vec{n})
      && \text{\jv{CollapsibleRowLayout}} \\
    \mid\;\; & \ccolp(\mu, a_{\mathrm{h}}, A, s, \vec{n})
      && \text{\jv{CollapsibleColumnLayout}} \\
    \mid\;\; & \fflow(\mu, A, a_{\mathrm{v}}, s, k_{\mathrm{items}}, k_{\mathrm{lines}}, \vec{n})
      && \text{\jv{FlowLayout}}
  \end{aligned}
\]
where $\vec{n}$ is a list of nodes, $s \in \Sc$ is the \jv{spacedBy} gap,
$\theta = (\sigma_0, s_0, a, s_{\min}, s_{\max})$ is a text specification (\jv{TextSpec}),
$a_{\mathrm{h}}$ and $a_{\mathrm{v}}$ are cross-axis alignments, $A$ is a
main-axis arrangement, and $k_{\mathrm{items}}, k_{\mathrm{lines}} \in \mathbb{N}$
are the flow's \jv{maxColumns} and \jv{maxLines}.
\end{definition}

A collapsible manager has the same shape as the plain one it extends, and a flow
has a row's shape plus its two limits. What distinguishes them is entirely in the
hooks of Section~\ref{sec:crow} and Section~\ref{sec:flow}.

\begin{definition}[Alignments]
\[
  a_{\mathrm{h}} \in \{\Start, \Center, \End\}, \qquad
  a_{\mathrm{v}} \in \{\Top, \Center, \Bottom\},
\]
\[
  A \in \{\Start, \Center, \End, \SBetween, \SEvenly, \SAround\}.
\]
\end{definition}

\begin{warning}[The defaults are not uniform]
\label{sec:defaults}
A document that omits an alignment does not get $\Start$ / $\Top$. The parser's
per-kind defaults (\jv{DefaultComponentParsers}, lines 55--113) are
\begin{center}
\small
\begin{tabular}{@{}ll@{}}
\toprule
\textbf{Kind} & \textbf{Default $(a_{\mathrm{h}}\text{ or }A,\; a_{\mathrm{v}}\text{ or }A)$} \\
\midrule
$\bbox$ with children & $(\Start, \Top)$ \\
$\bbox$ without children & $(\Center, \Center)$ \\
$\rrow$, $\ccol$ & $(\Start, \Top)$ \\
$\crow$, $\ccolp$ & $(\Center, \Center)$ \\
$\fflow$ & $(\Start, \Top)$, $k_{\mathrm{items}} = k_{\mathrm{lines}} = 2147483647$ \\
\bottomrule
\end{tabular}
\end{center}
and when both are present, \jv{horizontalAlignment} takes precedence over
\jv{horizontalArrangement} (\jv{RemoteComposeJsonParser}, lines 550--560).
Getting a collapsible manager's default wrong shifts every child.
\end{warning}

All six managers agree on the integer encoding written to the wire buffer:
\[
  \begin{gathered}
    \Start = 1,\quad \Center = 2,\quad \End = 3,\quad \Top = 4,\quad \Bottom = 5, \\
    \SBetween = 6,\quad \SEvenly = 7,\quad \SAround = 8 .
  \end{gathered}
\]

\begin{definition}[Spaced arrangements]
$\mathrm{spaced}(A)$ holds iff
\[
  A \in \{\SBetween,\; \SEvenly,\; \SAround\}.
\]
\end{definition}

\subsection{Output}

\begin{definition}[Layout tree]
A layout tree is
$\tau ::= \langle r, p, g, \vec{\tau}\,\rangle$ with $r$ a rectangle, $p$ the
node's padding, $g \in \mathbb{B}$ its visibility flag, and $\vec{\tau}$ its
laid-out children. We write $\tau.\gone$ for $g$, and call the node
\emph{visible} when $\neg g$.
\end{definition}

\begin{remark}[Why one bit suffices]
\jv{Component} carries six visibility values --- a base
$\jv{GONE} = 0$, $\jv{VISIBLE} = 1$, $\jv{INVISIBLE} = 2$ and the overrides
$\jv{OVERRIDE\_GONE} = 16$, $\jv{OVERRIDE\_VISIBLE} = 32$,
$\jv{OVERRIDE\_INVISIBLE} = 64$ --- read by
\[
  \mathrm{isGone}(v) = \begin{cases}
    (v \mathbin{\&} 16) = 16 & (v \gg 4) > 0 \\
    v = 0 & \text{otherwise,}
  \end{cases}
\]
so that an override, if present, wins outright. \jv{INVISIBLE} differs from
\jv{VISIBLE} only in painting, which is out of scope, and the \jv{visibility}
modifier that would let a document set the base value directly is also out of
scope. Everything that remains is $\mathrm{isGone}$, and one bit records it.
The distinction between a base value and an override does still matter, because
it decides whether a decision survives to the next pass; that shows up in
Caveat~\ref{sec:flowsticky}.
\end{remark}

\begin{warning}[Coordinate convention]
\label{sec:coords}
The $x$ and $y$ of a child are relative to the \emph{content origin} of its
parent: the parent's padding is \textbf{not} folded in. This matches
\jv{ComponentMeasure.setX} / \jv{setY}, which the managers fill in with offsets
measured inside the content box $(\,W - p_{\mathrm{h}}\,) \times (\,H -
p_{\mathrm{v}}\,)$. Window coordinates are recovered by
\jv{LayoutComponent.getLocationInWindow}, which adds each ancestor's padding
while walking up; the specification's $\mathrm{abs}$ function does the same.
The gold corpus records content-relative coordinates.
\end{warning}

An invisible child is still a child: it keeps its place in $\vec{\tau}$, it is
assigned coordinates, and the gold corpus records them
(Caveat~\ref{sec:goneplaced}). What it does not do is take up room.

%%% ==================================================================
\section{Derived attributes of a modifier chain}
\label{sec:derived}

Five functions read a chain. Three of them scan it in order and stop early;
that asymmetry is the source of the most surprising behaviour in the whole
algorithm.

\subsection{Padding: every modifier contributes}

\begin{definition}[Accumulated padding]
\[
  \Pp(\varepsilon) = (0,0,0,0), \qquad
  \Pp(\pad(p)\cdot\mu') = p + \Pp(\mu'), \qquad
  \Pp(m\cdot\mu') = \Pp(\mu') \text{ otherwise,}
\]
with $+$ componentwise.
\end{definition}

This mirrors the \jv{updatePadding} accumulation in
\jv{LayoutComponent.updateModifiers}: \emph{every} \jv{PaddingModifierOperation}
in the chain, wherever it sits, adds to \jv{mPaddingLeft} and friends. $\Pp(\mu)$
is what insets the children.

\subsection{Size descriptors: the first modifier wins}

\begin{definition}[Selected dimensions and bounds]
\[
  \Dimw(\varepsilon) = \Wrap, \qquad
  \Dimw(\wdt(d,\iota)\cdot\mu') = d, \qquad
  \Dimw(m\cdot\mu') = \Dimw(\mu') \text{ otherwise,}
\]
and $\Iw$ likewise returns the $\iota$ of the first $\wdt$. $\Dimh$, $\Ih$ are
the vertical mirrors.
\end{definition}

The default of $\Wrap$ is the reference implementation's: a component with no
size modifier wraps its content on that axis. ``First wins'' is
\jv{LayoutComponent.updateModifiers}, which assigns \jv{mWidthModifier} only
under the guard \jv{\&\& mWidthModifier == null}.

\begin{remark}[Encoder-level de-duplication]
In practice a chain rarely contains two width modifiers, because
\jv{RecordingModifier.setWidthModifier} keeps at most one slot per component and
a later write \emph{updates the existing element in place}, preserving its
position in the chain. Thus \jv{Modifier.width(100).weight(1)} records a single
$\Weight$ modifier at the position where the $\Exact(100)$ used to be, and the
$100$ is discarded. This is a property of document construction, not of layout,
and is not modelled; but it explains why the ``first wins'' rule is hard to
observe from the authoring API.
\end{remark}

\subsection{Modifier-defined size: the scan stops at the first size modifier}
\label{sec:orderquirk}

\begin{definition}[Modifier-defined width]
$\Dw(\mu, b) = g(\mu, 0, 0)$ where
\[
  \begin{aligned}
    g(\varepsilon, s, e) &= s + e, \\
    g(\wdt(d,\iota)\cdot\mu', s, e) &= s + \bigl(\base_b(d) \clampL \iota^-\bigr) + e,
      &&\text{(stop)} \\
    g(\pad(p)\cdot\mu', s, e) &= g(\mu', s + p_\ell, e + p_r), \\
    g(\hgt(\_)\cdot\mu', s, e) &= g(\mu', s, e).
  \end{aligned}
\]
$\Dh$ is the vertical mirror, accumulating $p_t$ and $p_b$ and stopping at the
first $\hgt$.
\end{definition}

This is \jv{LayoutComponent.computeModifierDefinedWidth}, which walks the same
list as \jv{updateModifiers} but \jv{break}s at the first width modifier. Two
consequences follow, and both are load-bearing:

\begin{enumerate}[leftmargin=1.6em]
  \item \textbf{Only padding declared \emph{before} the size modifier is added
    to the modifier-defined width.} Hence
    \[
      \Dw(\pad(8)\cdot\wdt(\Exact(50)), \max) = 66,
      \qquad
      \Dw(\wdt(\Exact(50))\cdot\pad(8), \max) = 50,
    \]
    while $\Pp$ --- and therefore the inset applied to the children --- is $8$ on
    every edge in both cases. \jv{Modifier.padding(8).size(50)} and
    \jv{Modifier.size(50).padding(8)} genuinely lay out differently.
  \item \textbf{A height modifier does not stop the width walk}, and vice versa.
    $\Dw(\pad(8)\cdot\hgt(\Exact(20))\cdot\wdt(\Exact(50)), \max) = 66$.
\end{enumerate}

%%% ==================================================================
\section{Intrinsic sizes}
\label{sec:intrinsic}

A second recursion, entirely independent of the first, computes what a node
would like to be if nothing constrained it. It is \jv{Component.minIntrinsicWidth}
and \jv{minIntrinsicHeight} together with the managers' overrides. Nothing in it
consults a measurement, so it is a function of the tree alone and terminates on
its own; \jv{FlowLayout} is the only consumer in scope, and it uses $\mih$ to
height the rows it has cut.

Throughout, the modifier-defined size is taken with $b = \min$, so that a $\Fill$
axis contributes $0$ rather than $\Om$ --- this is the \jv{isMin} argument to
\jv{computeModifierDefinedWidth}.

\begin{definition}[Intrinsic width]
\[
  \miw(n) = \begin{cases}
    \max\bigl(\Dw(\mu,\min),\; \sigma_w\bigr)
      & n = \leaf(\mu, \sigma) \\[2pt]
    \max\bigl(\Dw(\mu,\min),\; \max_i \miw(n_i)\bigr)
      & n = \bbox,\ \ccol \\[2pt]
    \max\bigl(\Dw(\mu,\min),\; \sum_i \miw(n_i)\bigr)
      & n = \rrow,\ \fflow \\[2pt]
    \Dw(\mu,\min) + \miw\bigl(\ell_{\Horz}(\vec{n})\bigr)
      & n = \crow \\[2pt]
    \Dw(\mu,\min) + \miw\bigl(\ell_{\Vrt}(\vec{n})\bigr)
      & n = \ccolp
  \end{cases}
\]
\end{definition}

\begin{definition}[Intrinsic height]
\[
  \mih(n) = \begin{cases}
    \max\bigl(\Dh(\mu,\min),\; \sigma_h\bigr)
      & n = \leaf(\mu, \sigma) \\[2pt]
    \max\bigl(\Dh(\mu,\min),\; \max_i \mih(n_i)\bigr)
      & n = \bbox,\ \rrow,\ \fflow \\[2pt]
    \max\bigl(\Dh(\mu,\min),\; \sum_i \mih(n_i)\bigr)
      & n = \ccol \\[2pt]
    \Dh(\mu,\min) + \mih\bigl(\ell_{\Horz}(\vec{n})\bigr)
      & n = \crow \\[2pt]
    \Dh(\mu,\min) + \mih\bigl(\ell_{\Vrt}(\vec{n})\bigr)
      & n = \ccolp
  \end{cases}
\]
\end{definition}

Empty maxima and empty sums are $0$, and the collapsible cases contribute $0$ for
the child term when the node has no children.

The pattern is that a manager sums along its main axis and maximises across it,
with two departures. \jv{BoxLayout} and \jv{ColumnLayout} override neither width
rule, so both inherit \jv{LayoutManager}'s maximum; and \jv{FlowLayout} inherits
\jv{RowLayout}'s width \emph{sum}, even though it wraps, so a flow's intrinsic
width is the width it would need on a single line.

\begin{definition}[Last standing]
$\ell_o(\vec{n})$ is the child of greatest $P_o$ (Section~\ref{sec:crow}), ties
going to the earliest, and $\bo$ for an empty list. It is
\jv{CollapsiblePriority.findLastStanding}, and it names the child a collapsible
manager would keep longest --- so its intrinsic size is the smallest the manager
can be reduced to without disappearing altogether.
\end{definition}

\begin{remark}
\jv{findLastStanding} is guarded by \jv{Header.FEATURE\_PRIORITY\_FIX}, which
defaults to enabled. The legacy branch took the first child instead; it is not
modelled.
\end{remark}

%%% ==================================================================
\section{Axis resolution}
\label{sec:axis}

Both axes of a node are resolved by the same function, applied once
horizontally and once vertically. It is the transcription of policy lines
90--154, where the reference implementation carries the same four results in
local variables.

\begin{definition}[Axis resolution]
\[
  \Rr(d, f, p, \delta, \mu_0, c^-, c^+, \hat c)
  \;=\; (a,\; c^{-\prime},\; \hat c\,',\; \omega)
\]
is defined by cases:
\[
\begin{array}{@{}l@{\;}l@{}}
\textbf{(F)} \text{ if } f: &
  \begin{cases}
    (c^+,\; \hat c,\; \hat c,\; \textsf{false})
      & \text{if } \val(d) = \bo \text{ or } d = \Exact, \\[2pt]
    (c^+\varphi,\; c^+\varphi - p,\; \hat c,\; \textsf{false})
      & \text{if } \val(d) = \varphi \in \Sc,\; d \neq \Exact,
  \end{cases} \\[16pt]
\textbf{(W)} \text{ else if } d = \Weight: &
  \bigl(\max(\mu_0, \delta),\; c^-,\; \hat c,\; \textsf{false}\bigr), \\[6pt]
\textbf{(C)} \text{ otherwise:} &
  \bigl(a,\; c^-,\; \hat c\,',\; \omega\bigr) \text{ with }
  \begin{aligned}[t]
    a &= \min\bigl(\max(\mu_0, c^-),\, c^+\bigr), \\
    \omega &= [\,d = \Wrap\,], \\
    \hat c\,' &= \begin{cases} \hat c & \omega \\ a - p & \neg\omega. \end{cases}
  \end{aligned}
\end{array}
\]
\end{definition}

The arguments are, in order: the axis' dimension descriptor $d$; the
\emph{fill} flag $f$ (Section~\ref{sec:fillflag}); the total padding $p$ on this
axis; the uncapped modifier-defined size $\delta$; the pre-capped measurement
$\mu_0 = \min(c^+_{\text{incoming}}, \delta)$; the constraints $c^-, c^+$ after
the \jv{widthIn} rewrite; and $\hat c = c^+ - p$, the maximum available to the
children. The results are the running measured size $a$, a possibly rewritten
minimum, the inset maximum handed down, and whether this axis wraps.

Three points deserve comment.

\begin{description}
  \item[Case (F) with an $\Exact$ descriptor.] A row whose own width is exact
    but whose children carry weights enters case~(F) because of the fill flag,
    and then $\jv{isExact()}$ wins and the row fills its whole slot. The
    declared exact width is overridden.
  \item[Case (W) is provisional.] A weighted axis is only given a placeholder
    here. Its real size is imposed later by its parent, through the tight
    constraints of the weight-resolution sweep (Sections~\ref{sec:row}
    and~\ref{sec:column}).
  \item[Inset wrap.] The rewrite $\hat c\,' = a - p$ in case~(C) is the effect of
    \jv{shouldApplyInsetWrap}, which is true for the default policy: when an
    axis does \emph{not} wrap, its children are capped by the size the node
    actually resolved to, not by the maximum it was offered.
\end{description}

\subsection{The fill flag}
\label{sec:fillflag}

\begin{definition}[Fill flags]
\[
  F_{\mathrm{h}}(n) =
  \begin{cases}
    \Dimw(\mu) = \Fill(\_) \;\vee\; \mathrm{someHW}(\vec{n})
      & n = \rrow(\mu, \_, \_, \_, \vec{n}) \\
    \Dimw(\mu) = \Fill(\_) & \text{otherwise,}
  \end{cases}
\]
\[
  F_{\mathrm{v}}(n) =
  \begin{cases}
    \Dimh(\mu) = \Fill(\_) \;\vee\; \mathrm{someVW}(\vec{n})
      & n = \ccol(\mu, \_, \_, \_, \vec{n}) \\
    \Dimh(\mu) = \Fill(\_) & \text{otherwise,}
  \end{cases}
\]
where $\mathrm{someHW}(\vec{n})$ holds iff some child has
$\Dimw = \Weight(\_)$, and $\mathrm{someVW}$ is its vertical mirror.
\end{definition}

A row containing weighted children must behave as though it were
\jv{fillMaxWidth}, since it has leftover space to hand out; \jv{RowLayout}
overrides \jv{isInHorizontalFill} to say so, and \jv{ColumnLayout} overrides
\jv{isInVerticalFill} symmetrically.

\begin{remark}
There is a latent discrepancy in the reference implementation here:
\jv{childrenHaveHorizontalWeights} tests \jv{instanceof LayoutManager} while
\jv{computeWrapSize} tests \jv{instanceof LayoutComponent}. The two predicates
differ for non-manager layout components. No document in the corpus
distinguishes them, and the specification does not model the difference.
\end{remark}

\subsection{Weight totals}
\label{sec:weights}

\[
  \Theta_{\mathrm{h}}(\vec{n}) = \!\!\sum_{i \,:\, \neg \tau_i.\gone}\!\! \wt(\Dimw(\mu_i)),
  \qquad
  \Theta_{\mathrm{v}}(\vec{n}) = \!\!\sum_{i \,:\, \neg \tau_i.\gone}\!\! \wt(\Dimh(\mu_i)).
\]

The sums --- and the ``some child is weighted'' tests that gate them --- range
over the \emph{visible} children only. \jv{RowLayout} guards both
\jv{hasWeights} and \jv{totalWeights} with \jv{isGone} (lines 379 and 399), so a
child that a collapsing sweep or a flow segmentation has dropped neither claims
a share of the leftover space nor dilutes anybody else's. For the three
managers that cannot make a child invisible this is vacuous; for the other
three it is load-bearing.

\begin{warning}[A bare \texttt{weight} is axis-relative]
\label{sec:axisweight}
This is a fact about authoring rather than about layout, but it decides which of
$\Dimw$ and $\Dimh$ a document's \jv{"weight"} lands in, and so which of the two
sums above sees it. \jv{DefaultModifierParsers} (lines 145--155) consults
\jv{parser.isParentVertical()} and routes the key to \jv{verticalWeight} inside
a \jv{column} or \jv{collapsibleColumn} and to \jv{horizontalWeight} everywhere
else --- note that \jv{flow} counts as \emph{horizontal}, being row-major.
\jv{"collapsiblepriority"} (lines 193--216) resolves its orientation the same
way. The explicit \jv{"horizontalweight"} and \jv{"verticalweight"} keys bypass
the rule.
\end{warning}

%%% ==================================================================
\section{The measure judgement}
\label{sec:measure}

\begin{definition}[Measure]
$\Mm(n, \kappa) = \tau$ is defined by the following sequence, where
$\kappa = (w^-, w^+, h^-, h^+)$ and $\mu$ is $n$'s chain.
\end{definition}

\paragraph{Step 1 --- padding and modifier-defined sizes (policy 57--60).}
\[
  p = \Pp(\mu), \qquad
  \delta_w = \Dw(\mu, \max), \qquad
  \delta_h = \Dh(\mu, \max),
\]
\[
  \mu^0_w = \min(w^+, \delta_w), \qquad
  \mu^0_h = \min(h^+, \delta_h).
\]
For a $\Fill$ axis $\delta = \Om$ and the $\min$ collapses to $w^+$; this is
precisely why $\Om$ must be a real number.

\paragraph{Step 2 --- \texttt{widthIn} / \texttt{heightIn} rewrite (policy 70--79).}
\[
  w^-_1 = w^- \clampL \Iw^-, \quad
  w^+_1 = w^+ \clampH \Iw^+, \quad
  h^-_1 = h^- \clampL \Ih^-, \quad
  h^+_1 = h^+ \clampH \Ih^+ .
\]

\paragraph{Step 3 --- inset maxima (policy 81--82).}
\[
  \hat c_w = w^+_1 - p_{\mathrm{h}}, \qquad
  \hat c_h = h^+_1 - p_{\mathrm{v}} .
\]

\paragraph{Step 4 --- resolve both axes (policy 90--154).}
\[
  (a_w, w^-_2, \hat w, \omega_w)
    = \Rr\bigl(\Dimw(\mu),\, F_{\mathrm{h}}(n),\, p_{\mathrm{h}},\, \delta_w,\,
               \mu^0_w,\, w^-_1,\, w^+_1,\, \hat c_w\bigr),
\]
\[
  (a_h, h^-_2, \hat h, \omega_h)
    = \Rr\bigl(\Dimh(\mu),\, F_{\mathrm{v}}(n),\, p_{\mathrm{v}},\, \delta_h,\,
               \mu^0_h,\, h^-_1,\, h^+_1,\, \hat c_h\bigr).
\]

\paragraph{Step 5 --- enforce constraints (policy 163--168).}
\[
  a^1_w = \min\bigl(\max(a_w, w^-_2),\, w^+_1\bigr), \qquad
  a^1_h = \min\bigl(\max(a_h, h^-_2),\, h^+_1\bigr).
\]

\paragraph{Step 6 --- fully determined axes (policy 170--175).}
\[
  a^2_w = \begin{cases} w^+_1 & w^-_2 = w^+_1 \\ a^1_w & \text{else,}\end{cases}
  \qquad
  a^2_h = \begin{cases} h^+_1 & h^-_2 = h^+_1 \\ a^1_h & \text{else.}\end{cases}
\]
A tight constraint wins outright, whatever the content wants. This is why a
$\Wrap$ root still fills its viewport: \jv{RootLayoutComponent} measures the
document tightly.

\paragraph{Step 7 --- measure the children (policy 181--336).}
Exactly one of two hooks runs.

\emph{Wrap branch}, taken when $\omega_w \vee \omega_h$ (policy 182--205). Let
$\kappa_{\mathrm{wrap}} = (w^-_2, \hat w, h^-_2, \hat h)$ and let
$(\varsigma, \vec{\tau}\,)$ be the manager's wrap hook applied to
$(\vec{n}, \kappa_{\mathrm{wrap}})$, returning a content size $\varsigma$ and the
measured children. Then
\[
  a^3_w = \begin{cases}
    \max(\varsigma_w + p_{\mathrm{h}},\, w^-_2) & \omega_w \\
    a^2_w & \neg\omega_w
  \end{cases}
  \qquad
  a^3_h = \begin{cases}
    \max(\varsigma_h + p_{\mathrm{v}},\, h^-_2) & \omega_h \\
    a^2_h & \neg\omega_h .
  \end{cases}
\]

\emph{Size branch}, otherwise (policy 324--335). With
$\kappa_{\mathrm{size}} = (0,\; a^2_w - p_{\mathrm{h}},\; 0,\; a^2_h -
p_{\mathrm{v}})$, the manager's size hook produces $\vec{\tau}$, and
$a^3_w = a^2_w$, $a^3_h = a^2_h$. The children see the node's resolved content
box.

\paragraph{Step 8 --- final size (policy 349--357).}
\[
  W = \bigl(\min(\max(a^3_w, w^-_2),\, w^+_1)\bigr) \clampL \Iw^- \clampH \Iw^+,
\]
\[
  H = \bigl(\min(\max(a^3_h, h^-_2),\, h^+_1)\bigr) \clampL \Ih^- \clampH \Ih^+ .
\]

\paragraph{Step 9 --- placement (policy 364).}
The content box is
\[
  W_{\!*} = W - p_{\mathrm{h}}, \qquad H_{\!*} = H - p_{\mathrm{v}},
\]
and the manager's placement hook assigns each measured child an offset inside
it, returning the placed children together with a flag $g$ saying whether the
manager has decided it is itself invisible:
\[
  \Mm(n,\kappa) \;=\; \bigl\langle (0, 0, W, H),\; p,\; g,\; \mathrm{place}(\ldots) \bigr\rangle .
\]
The node's own $x$ and $y$ are left at $0$; its parent sets them. Only the two
collapsible managers can return $g = \textsf{true}$, and only when nothing
survived their final sweep (Section~\ref{sec:crow}); for every other kind
$g = \textsf{false}$.

\begin{definition}[Document layout]
\[
  \mathrm{layout}(n, W, H) = \Mm\bigl(n,\, \kappa_{\mathrm{tight}}(W, H)\bigr).
\]
\end{definition}

\begin{definition}[Absolute coordinates]
$\mathrm{abs}(\tau, o_x, o_y)$ lists every rectangle of $\tau$ in pre-order,
translated into window coordinates:
\[
  \mathrm{abs}\bigl(\langle r, p, g, \vec\tau\,\rangle, o_x, o_y\bigr)
    = (o_x + r.x,\, o_y + r.y,\, r.w,\, r.h) \;::\;
      \bigl[\,\mathrm{abs}(\tau_i,\; o_x + r.x + p_\ell,\; o_y + r.y + p_t)\,\bigr]_i .
\]
The parallel traversal that lists the flags $g$ in the same order is written
$\mathrm{abs}_{\gone}$; the conformance harness of
Section~\ref{sec:conformance} checks both.
\end{definition}

%%% ==================================================================
\section{Per-manager rules}

Six managers are specified, and they differ in exactly three places: the wrap
pass, the size pass, and the placement pass. Everything in
Section~\ref{sec:measure} is shared. Two of the six insert an extra structural
step of their own before those hooks run --- the collapsible managers decide
which children to keep, and the flow manager decides which row each child joins
--- but both steps are pure functions of children that have already been
measured, so they sit inside the hooks rather than beside them.

\subsection{Box}

\begin{definition}[Box hooks]
\[
  \mathrm{wrap}_{\bbox}(\vec{n}, \kappa) =
    \Bigl(\bigl(\textstyle\max_i \tau_i.w,\; \max_i \tau_i.h\bigr),\; \vec{\tau}\,\Bigr),
  \qquad \tau_i = \Mm\bigl(n_i,\, \kappa_{\mathrm{loose}}(\kappa.w^+, \kappa.h^+)\bigr),
\]
\[
  \mathrm{size}_{\bbox}(\vec{n}, \kappa) = \bigl[\,\Mm(n_i, \kappa)\,\bigr]_i ,
\]
\[
  \mathrm{place}_{\bbox}(\vec{\tau}, a_{\mathrm{h}}, a_{\mathrm{v}}, W_{\!*}, H_{\!*})
    = \bigl[\,\tau_i \text{ at } \bigl(X(a_{\mathrm{h}}, W_{\!*}, \tau_i.w),\;
                                       Y(a_{\mathrm{v}}, H_{\!*}, \tau_i.h)\bigr)\,\bigr]_i ,
\]
with $X$, $Y$ from Section~\ref{sec:cross}, and the empty maxima taken to be $0$.
\end{definition}

Every child of a box sees identical constraints, and the box takes the pointwise
maximum. Children overlap; they are painted in list order.

\subsection{Row}
\label{sec:row}

\paragraph{Wrap pass.}
Let $k = |\vec{n}|$. If no child is weighted, the row performs a single
left-to-right sweep in which each child sees the space its predecessors did not
take:
\[
  \tau_i = \Mm\bigl(n_i,\, (0,\, \rho_{i-1},\, 0,\, \kappa.h^+)\bigr),
  \qquad \rho_0 = \kappa.w^+, \quad \rho_i = \rho_{i-1} - \tau_i.w .
\]
If some child is weighted, the sweep is split. First the unweighted children are
measured to find the space left over,
\[
  \rho^\star = \kappa.w^+ - \!\!\sum_{i \,:\, \Dimw(\mu_i) \neq \Weight(\_)}\!\! \tau_i.w ,
\]
then the measuring sweep runs, pinning each weighted child to its share:
\[
  \tau_i = \begin{cases}
    \Mm\bigl(n_i,\, \kappa_{\mathrm{tight}}(c_i, \cdot)\bigr),\;
      c_i = \dfrac{\wt(\Dimw(\mu_i)) \cdot \rho^\star}{\Theta_{\mathrm{h}}(\vec{n})}
      & \Dimw(\mu_i) = \Weight(\_) \\[10pt]
    \Mm\bigl(n_i,\, (0,\, \rho_{i-1},\, 0,\, \kappa.h^+)\bigr) & \text{otherwise,}
  \end{cases}
\]
where the running budget $\rho$ advances only on unweighted children. (Weighted
children receive the vertical constraint $[0, \kappa.h^+]$.) The content size is
\[
  \varsigma_w = \sum_i \tau_i.w + s\,(k-1), \qquad
  \varsigma_h = \max_i \tau_i.h ,
\]
with the spacing term taken to be $0$ when $k = 0$.

\begin{remark}
Because $\Mm$ is a function, re-measuring an unweighted child in the second
sweep yields exactly the value the first sweep computed for it. The reference
implementation's two physical passes therefore collapse into one ordered sweep
without changing the result.
\end{remark}

\paragraph{Size pass.}
A plain shrinking sweep that forwards the minimum unchanged and shrinks only the
maximum:
\[
  \tau_i = \Mm\bigl(n_i,\, (\kappa.w^-,\, \rho_{i-1},\, \kappa.h^-,\, \kappa.h^+)\bigr),
  \qquad \rho_0 = \kappa.w^+, \quad \rho_i = \rho_{i-1} - \tau_i.w .
\]
Weights are \emph{not} resolved here; they are deferred to placement.

\paragraph{Placement.}
If some child is weighted, each weighted child is re-measured tightly against
its share of the space left inside the \emph{final} content width, preserving
the height it already measured:
\[
  \mathrm{avail} = W_{\!*} - \!\!\sum_{i \,:\, \Dimw(\mu_i) \neq \Weight(\_)}\!\! \tau_i.w,
  \qquad
  \tau_i' = \Mm\Bigl(n_i,\; \kappa_{\mathrm{tight}}\bigl(
      \tfrac{\wt(\Dimw(\mu_i)) \cdot \mathrm{avail}}{\Theta_{\mathrm{h}}(\vec{n})}
      \clampL \Iw^- \clampH \Iw^+ ,\;\; \tau_i.h \bigr)\Bigr).
\]
Unweighted children keep $\tau_i' = \tau_i$. Then, with
\[
  T = \sum_i \tau_i'.w, \qquad
  C = T + s\,(k-1),
\]
the arrangement $\Aa$ of Section~\ref{sec:arrange} yields an initial offset
$t_0$ and an inter-child gap $g$, and the children are laid out by
\[
  x_1 = t_0, \qquad
  x_{i+1} = \begin{cases}
    x_i & \tau_i' \text{ is } \gone \\
    x_i + \tau_i'.w + [\,\mathrm{spaced}(A)\,]\cdot g + s & \text{otherwise,}
  \end{cases}
  \qquad
  y_i = Y(a_{\mathrm{v}}, H_{\!*}, \tau_i'.h).
\]
The sums $T$ and $C$, and the count $k$ that $\Aa$ divides by, likewise range
over the visible children only; $\Theta_{\mathrm{h}}$ skips the invisible ones
too, so a collapsed child neither claims a share of the leftover space nor
dilutes anyone else's.

\begin{warning}[A collapsed child is still positioned]
\label{sec:goneplaced}
Only the \emph{cursor advance} is skipped, not the assignment. \jv{setX} and
\jv{setY} run at lines 546--547, before the \jv{isGone} test at line 548, so an
invisible child is recorded at exactly the coordinates its next visible sibling
will occupy. Nothing is painted there, but the gold dumps record the rectangle
and the specification reproduces it. This matters only for the collapsible and
flow managers, which are the only ones in scope that can make a child invisible.
\end{warning}

\begin{remark}[The weight loop]
The reference implementation wraps weight resolution in
\jv{while (checkWeights)}, whose body re-runs only when \jv{applyVisibility}
returns true. \jv{LayoutManager.applyVisibility} returns \jv{false}
unconditionally (lines 76--79), and \textbf{no} subclass overrides it --- not
even the collapsible managers, which do their visibility work in
\jv{computeVisibleChildren} instead. The loop is therefore always a single pass,
for every manager, and is modelled as straight-line code.
\end{remark}

\subsection{Column}
\label{sec:column}

The column is the vertical mirror of the row in its wrap pass and its
placement pass. It is \textbf{not} a mirror in its size pass.

\begin{warning}[Row/Column asymmetry]
\label{sec:rowcolasym}
\jv{ColumnLayout.computeSize} is guarded by
\begin{center}
\jv{private static final boolean DIRECT\_WEIGHT\_CALCULATION = true;}
\end{center}
and resolves vertical weights \emph{during the size pass}, in two sweeps with
different running budgets. \jv{RowLayout.computeSize} does no such thing.
Specifying the column as a mirror of the row disagrees with the gold corpus.
\end{warning}

\paragraph{Size pass with weights.}
When no child is weighted, the column runs the vertical mirror of the row's
shrinking sweep. When some child is weighted:

\emph{First sweep.} The unweighted children are measured against their own
running budget and their total height recorded:
\[
  N = \!\!\sum_{i \,:\, \Dimh(\mu_i) \neq \Weight(\_)}\!\! \tau_i.h,
  \qquad
  \tau_i = \Mm\bigl(n_i,\, (\kappa.w^-,\, \kappa.w^+,\, \kappa.h^-,\, \eta_{i-1})\bigr),
\]
with $\eta_0 = \kappa.h^+$ and $\eta$ decreasing only on unweighted children.

\emph{Second sweep.} With $\mathrm{avail} = \kappa.h^+ - N$, every child is
measured again, weighted ones tightly:
\[
  \tau_i' = \begin{cases}
    \Mm\bigl(n_i,\, (\kappa.w^-,\, \kappa.w^+,\, c_i,\, c_i)\bigr),\;
      c_i = \dfrac{\mathrm{avail}\cdot\wt(\Dimh(\mu_i))}{\Theta_{\mathrm{v}}(\vec{n})}
      & \Dimh(\mu_i) = \Weight(\_) \\[10pt]
    \Mm\bigl(n_i,\, (\kappa.w^-,\, \kappa.w^+,\, \kappa.h^-,\, \eta'_{i-1})\bigr)
      & \text{otherwise,}
  \end{cases}
\]
where the second budget $\eta'$ now shrinks by \emph{every} child, weighted ones
included. An unweighted child can therefore measure differently in the two
sweeps; it is the second result that is kept.

\paragraph{Wrap and placement.}
The vertical mirrors of the row's, with $\Dimh$, $\Theta_{\mathrm{v}}$,
$\varsigma_h = \sum_i \tau_i.h + s(k-1)$, $\varsigma_w = \max_i \tau_i.w$, and
the cross-axis offset $X(a_{\mathrm{h}}, W_{\!*}, \tau_i.w)$.

\subsection{Collapsible row}
\label{sec:crow}

A \jv{CollapsibleRowLayout} is a \jv{RowLayout} that first decides which of its
children it can afford to keep, and then defers to \jv{RowLayout} for the
arrangement. The decision --- \jv{computeVisibleChildren}, lines 239--326 --- is
taken in two phases, and it is retaken from scratch at each of the manager's
three entry points.

\paragraph{Priorities.}

\begin{definition}[Collapsible priority]
For an orientation $o \in \{\Horz, \Vrt\}$,
\[
  P_o(n) = \begin{cases}
    v & \text{the chain of } n \text{ contains } \pri(o, v), \\
    \Om & \text{otherwise.}
  \end{cases}
\]
\end{definition}

A child that declares no priority --- or declares one for the \emph{other} axis
--- therefore holds the largest priority there is, and is the last to be dropped.

\begin{definition}[Priority order]
$a \prec_o b$ iff $P_o(a) \geq P_o(b) + 1$. Write $\mathrm{sort}_o$ for the
stable sort by $\prec_o$: highest priority first, so that the tail of the sorted
list is what goes first.
\end{definition}

The reference comparator is \jv{(int) (priority2 - priority1)}, and the $+1$
encodes its truncation towards zero: an \jv{int} result is negative only for a
difference of at least one, so two priorities less than a unit apart compare
equal and, the sort being stable, keep their declaration order.

\begin{warning}[Fractional priorities are unspecified]
Within one unit the truncating comparator is not a strict weak ordering, and the
reference implementation's output then depends on the internals of Java's
\jv{TimSort}. Integer priorities --- everything the corpus uses --- are
unaffected, and the specification commits to nothing else.
\end{warning}

\paragraph{Phase 1: the measuring sweep (lines 251--281).}
Each child is measured once, in declaration order, against a budget that shrinks
by what its predecessors took --- but the budget is only \emph{offered} to a
child that is itself a collapsible row or that carries a horizontal weight:
\[
  \tau_i = \Mm\bigl(n_i,\; (0,\; \beta^{\ast}_i,\; 0,\; \kappa.h^+)\bigr),
  \qquad
  \beta^{\ast}_i = \begin{cases}
    \beta_{i-1} & n_i = \crow(\ldots) \text{ or } \Dimw(\mu_i) = \Weight(\_) \\
    \Om & \text{otherwise,}
  \end{cases}
\]
\[
  \beta_0 = \kappa.w^+,
  \qquad
  \beta_i = \begin{cases}
    \beta_{i-1} & \tau_i \text{ is } \gone \\
    \beta_{i-1} - \tau_i.w & \text{otherwise.}
  \end{cases}
\]

\begin{remark}
The $\Om$ is the point of the phase. A child measured against the space that is
actually left would report a width that fits by construction, and phase~2 would
have nothing to discover. Measuring it against \jv{Float.MAX\_VALUE} instead
makes it report the width it \emph{wants}, which is what the fit test needs.
\end{remark}

\paragraph{Phase 2: the collapsing sweep (lines 290--325).}
This phase re-enters nothing: it works purely on the sizes phase~1 established.

\begin{definition}[Collapsing sweep]
Let $e(\tau) = \tau.w$ for a row and $\tau.h$ for a column. For a budget $B$,
\[
  \Sigma(B, \varepsilon, \alpha, f) = \varepsilon,
\]
\[
  \Sigma\bigl(B,\; (i,\tau)\cdot\vec{x},\; \alpha,\; f\bigr) = \begin{cases}
    (i, \textsf{true}) \cdot \Sigma(B, \vec{x}, \alpha, f)
      & f \vee \tau.\gone, \\[2pt]
    (i, \textsf{true}) \cdot \Sigma(B, \vec{x}, \alpha, \textsf{true})
      & \alpha + e(\tau) > B, \\[2pt]
    (i, \textsf{false}) \cdot \Sigma\bigl(B, \vec{x}, \alpha + e(\tau), \textsf{false}\bigr)
      & \text{otherwise.}
  \end{cases}
\]
\end{definition}

Three things are load-bearing. The refusal flag $f$ is \textbf{sticky}: once one
child has been refused, so is every child after it in priority order, however
small. A child that arrived already $\gone$ is skipped \emph{without} setting the
flag, so it costs its successors nothing. And $\alpha$ accumulates over the
survivors only.

\begin{definition}[Collapse decision]
$G_o(B, \vec{n}, \vec{\tau}\,)$ is the list of flags produced by
$\Sigma(B, \Pi, 0, \textsf{false})$, read back into declaration order, where
$\Pi$ is the indexed child list --- sorted by $\mathrm{sort}_o$ when at least one
child carries a $\pri$ modifier, and left in declaration order otherwise.
\end{definition}

\begin{warning}[The sort gate ignores orientation]
The \jv{hasPriorities} flag is set by the mere presence of a
\jv{CollapsiblePriorityModifier}, whatever axis it names. A child carrying a
\emph{vertical} priority is therefore enough to make a \emph{row} sort its
children --- and $P_\Horz$ then returns $\Om$ for that child, so it sorts to the
front and collapses last.
\end{warning}

\paragraph{The three hooks.}
Write $\vec{\tau} = \Phi(\vec{n}, \kappa)$ for the result of phase~1 and $\nu$
for the number of children it left visible.

\emph{Wrap pass.}
\[
  \varsigma_w = \begin{cases}
    \min\bigl(\kappa.w^+,\; \sum_{i \notin G} \tau_i.w\bigr) & \omega_w \\[4pt]
    \sum_i \tau_i.w + s\,(\nu - 1) & \neg\omega_w
  \end{cases}
  \qquad
  \varsigma_h = \max_i \tau_i.h,
\]
with $G = G_\Horz(\kappa.w^+, \vec{n}, \vec{\tau}\,)$ and the spacing term taken
as $0$ for an empty child list. The two branches disagree about more than the
clamp: the wrapping one sums the \emph{survivors} and contains no spacing at all,
the other sums \emph{everything} phase~1 measured and does. The height is the
tallest child of phase~1 either way, collapsed or not.

\emph{Size pass.} Phase~1 alone. Whatever $G$ this pass would compute is
discarded, because the placement pass clears every visibility override and takes
the decision again.

\emph{Placement.} The decision is taken a final time, and it is this one that is
observable:
\[
  \vec{\tau}\,' = \text{apply } G_\Horz(W, \vec{n}, \vec{\tau}\,),
  \qquad
  \mathrm{place}_{\crow} = \mathrm{place}_{\rrow}\bigl(\vec{\tau}\,', A, a_{\mathrm{v}}, s, W_{\!*}, H_{\!*}\bigr),
\]
and the manager declares \emph{itself} $\gone$ iff no $\tau'_i$ survives.

\begin{warning}[The budget includes the padding]
The final decision is taken against $W$, the manager's full measured width, and
\textbf{not} against $W_{\!*} = W - p_{\mathrm{h}}$, the content box the children
are then placed inside. A padded collapsible row admits children its content box
cannot hold, and they overflow it by up to $p_{\mathrm{h}}$.
\end{warning}

\begin{warning}[The fit test is blind to \texttt{spacedBy}]
\label{sec:collapsespaced}
$\Sigma$ compares $\alpha + e(\tau)$ against $B$ with no spacing term, and the
placement loop then adds $s$ between every surviving pair anyway. A spaced
collapsible layout therefore admits one child too many and overflows by up to
$s\,(\nu - 1)$; see Appendix~\ref{app:examples}, Example~6.
\end{warning}

\begin{remark}[Why the decision is retaken]
Collapsible visibility is written as an \jv{OVERRIDE\_GONE} bit rather than a
base value, and every entry point begins by calling
\jv{clearVisibilityOverride}. Only the last decision --- the one placement takes
--- reaches the output. An empty collapsible manager declares itself $\gone$
vacuously, having no visible child.
\end{remark}

\subsection{Collapsible column}
\label{sec:ccol}

The vertical mirror, and this time genuinely a mirror: unlike \jv{ColumnLayout},
whose size pass departs from \jv{RowLayout}'s (Caveat~\ref{sec:rowcolasym}),
\jv{CollapsibleColumnLayout} is a line-for-line transcription of
\jv{CollapsibleRowLayout} with the axes exchanged. Read
Section~\ref{sec:crow} with $\Dimh$ for $\Dimw$, $\Vrt$ for $\Horz$, $e(\tau) =
\tau.h$, the budget $H$ for $W$, and $\ccolp$ for $\crow$ in the phase-1 test.

Note in particular that the collapsible column does \emph{not} inherit
\jv{ColumnLayout}'s \jv{DIRECT\_WEIGHT\_CALCULATION} asymmetry, because it
overrides \jv{computeSize} itself.

\subsection{Flow}
\label{sec:flow}

A \jv{FlowLayout} is a \jv{RowLayout} that cuts its children into rows and then
runs the inherited row machinery once per row. Segmentation is a pure function
of the measured children, like the collapsing sweep, and is likewise re-run by
each entry point.

\paragraph{Extents.}

\begin{definition}[Segmentation extents]
\[
  e_i = \begin{cases}
    (0, 0) & \tau_i \text{ is } \gone, \\
    \bigl(\Iw^-(\mu_i) \text{ or } 0,\;\; \mih(n_i)\bigr)
      & \Dimw(\mu_i) = \Weight(\_), \\
    (\tau_i.w,\; \tau_i.h) & \text{otherwise.}
  \end{cases}
\]
\end{definition}

A weighted child is not measured for segmentation at all: it is credited with its
\jv{widthIn} minimum, or nothing, and with its intrinsic height
(Section~\ref{sec:intrinsic}). Its real width is settled later, once the row it
landed in is known.

\paragraph{The segmentation walk (lines 225--277).}
The state is $(c, \theta, \Theta, \varrho, \nu, f)$: the current row's width and
height, the height already stacked, the number of rows opened, the number of
items in the current row, and the sticky refusal flag. It starts at
$(0, 0, 0, 1, 0, \textsf{false})$ --- note $\varrho = 1$ --- and the box is
$(\hat W, \hat H)$.

For each child in declaration order: if $f$, the child is refused. Otherwise let
\[
  \mathrm{brk} = \bigl(e_i.w + c > \hat W\bigr) \;\vee\; \bigl(\nu \geq k_{\mathrm{items}}\bigr)
\]
be the test that the current row cannot take it. If $\mathrm{brk}$ holds and no
further row may be opened,
\[
  \varrho \geq k_{\mathrm{lines}} \;\vee\; \Theta + \theta \geq \hat H,
\]
the child is refused and $f$ is set. Otherwise the child joins row
$\varrho' - 1$, and the state advances by
\[
  \begin{gathered}
    \varrho' = \varrho + [\mathrm{brk}], \qquad
    \Theta' = \Theta + [\mathrm{brk}]\cdot\theta, \\
    c' = [\neg\mathrm{brk}]\cdot c + e_i.w + s, \qquad
    \theta' = \max\bigl([\neg\mathrm{brk}]\cdot\theta,\; e_i.h\bigr), \qquad
    \nu' = [\neg\mathrm{brk}]\cdot\nu + 1,
  \end{gathered}
\]
every right-hand side reading the pre-update state. Write
$R = \{0, \ldots, |R| - 1\}$ for the rows opened, $|R| = \max\bigl(1, 1 + \max_i
r_i\bigr)$, and $i \in r$ for membership.

\begin{warning}[A flow opens a row before it starts]
$\varrho$ counts from $1$, because the reference implementation appends an empty
row before the walk. $k_{\mathrm{lines}}$ therefore bites one row earlier than
the name suggests. Both limits default to \jv{Integer.MAX\_VALUE} $=
2147483647$.
\end{warning}

\begin{warning}[\texttt{spacedBy} is added after the last item of a row]
$c'$ adds $s$ after \emph{every} child, including the one that ends a row, so a
spaced flow breaks one item earlier than the arithmetic suggests --- the phantom
trailing gap is charged against the next child's fit test.
\end{warning}

\begin{warning}[A refused child does not stay refused]
\label{sec:flowsticky}
\jv{segmentComponents} records a refusal with \jv{setVisibility(GONE)}, which
writes the component's \emph{base} visibility rather than an override, and
nothing ever clears it. The refusal therefore survives into the next pass, which
reads the child back through the first case of $e_i$ as a zero-width item, finds
that a zero-width item fits, and lets it rejoin a row --- at whatever position the
cursor had reached, which may be well outside the container. Appendix
\ref{app:examples}, Example~8 shows this happening; \file{flow\_max\_lines} is
tagged \jv{suspicious} in the gold suite for exactly it.
\end{warning}

\paragraph{Wrap pass (lines 290--331).}
Every child is measured against the full box, the children are segmented, and
then each row in turn is handed to the inherited \jv{RowLayout.computeWrapSize}
--- which re-measures that row's unweighted members against a budget shrinking
along the row, and its weighted members against their share of what is left,
exactly as in Section~\ref{sec:row}. With
\[
  \mathrm{rw}(r) = \!\!\sum_{i \in r,\; \text{visible}}\!\! \tau_i.w
    \;+\; s\,\bigl(|r|_{\vis} - 1\bigr)
  \quad (0 \text{ for an empty row}),
\]
the content size is
\[
  \varsigma_w = \min\Bigl(\max\bigl(\kappa.w^-,\; \max_{r} \mathrm{rw}(r)\bigr),\; \kappa.w^+\Bigr),
  \qquad
  \varsigma_h = \min\Bigl(\max\bigl(\kappa.h^-,\; \sum_{r} \max_{i \in r} \tau_i.h\bigr),\; \kappa.h^+\Bigr).
\]
This is the only wrap hook in the specification that clamps into its own incoming
constraints rather than leaving that to Step~8.

\paragraph{Size pass (lines 333--354).}
A single uniform measurement against the content box, followed by segmentation
and the refusals it decides. There is no per-row re-measurement here: the
shrinking budget that would provide one is commented out at line~350 of the
reference implementation, which makes the two nominal sweeps identical.

\paragraph{Placement (lines 356--399).}
Four steps.

\begin{enumerate}[leftmargin=1.6em]
  \item \emph{Re-measure.} Every child that is neither $\gone$ nor weighted is
    measured again against $(W_{\!*}, H_{\!*})$; the others keep the size they
    have. A $\gone$ child is thus left holding a stale size that segmentation
    will not read anyway.
  \item \emph{Re-segment} against $(W_{\!*}, H_{\!*})$, giving the final rows and
    the final refusals.
  \item \emph{Stack.} Each row is given the intrinsic height
    \[
      \eta(r) = \max\Bigl(\Dh(\mu, \min),\;\; \max_{i \in r} \mih(n_i)\Bigr),
    \]
    where $\mu$ is the \textbf{flow's own} chain --- the protected
    \jv{RowLayout.minIntrinsicHeight} is invoked on \jv{this}, not on the row.
    The rows begin at $y_0 = Y\bigl(a_{\mathrm{v}}, H_{\!*}, \sum_r \eta(r)\bigr)$.
  \item \emph{Lay out each row} as an ordinary row inside a box
    $W_{\!*} \times \eta(r)$, shift its children down by the running $y$, and
    advance
    \[
      y \;\mathrel{+}=\; \max_{i \in r,\; \text{visible}} \tau_i.h .
    \]
\end{enumerate}

\begin{warning}[Two different row heights]
Step~3 resolves the stack, and the cross-axis alignment \emph{inside} each row,
against the row's \emph{intrinsic} height $\eta(r)$; step~4 advances the cursor
by the row's \emph{measured} height. These are different quantities, and nothing
reconciles them. Where a row measures taller than it intrinsically is, the rows
gape apart and the last of them runs past the bottom of the flow; where it
measures shorter, they overlap.
\end{warning}

%%% ==================================================================
\section{Alignment and arrangement}

\subsection{Cross axis}
\label{sec:cross}

\begin{definition}[Cross-axis offsets]
\[
  X(a_{\mathrm{h}}, L, c) = \begin{cases}
    0 & a_{\mathrm{h}} = \Start \\
    (L - c)/2 & a_{\mathrm{h}} = \Center \\
    L - c & a_{\mathrm{h}} = \End
  \end{cases}
  \qquad
  Y(a_{\mathrm{v}}, L, c) = \begin{cases}
    0 & a_{\mathrm{v}} = \Top \\
    (L - c)/2 & a_{\mathrm{v}} = \Center \\
    L - c & a_{\mathrm{v}} = \Bottom .
  \end{cases}
\]
\end{definition}

$X$ is used for the cross axis of a column and for the horizontal axis of a box;
$Y$ for the cross axis of a row and the vertical axis of a box. Note that
neither clamps: a child larger than its container gets a negative offset, and
overflows symmetrically under $\Center$.

\subsection{Main axis}
\label{sec:arrange}

\begin{definition}[Arrangement]
Given the container's main extent $L$, the children's total extent \emph{with}
spacing $C$, their total extent \emph{without} spacing $T$, and their count $k$,
\[
  \Aa(A, L, C, T, k) = (t_0, g)
\]
where
\[
\begin{array}{lll}
  \Aa(\Start) &=& (0,\; 0) \\
  \Aa(\End) &=& (L - C,\; 0) \\
  \Aa(\Center) &=& \bigl((L - C)/2,\; 0\bigr) \\
  \Aa(\SBetween) &=& \begin{cases}
      \bigl(0,\; (L - T)/(k-1)\bigr) & k > 1 \\
      \bigl((L - C)/2,\; 0\bigr) & k \leq 1
    \end{cases} \\[10pt]
  \Aa(\SEvenly) &=& (g,\; g), \quad g = (L - T)/(k+1) \\[2pt]
  \Aa(\SAround) &=& (g/2,\; g), \quad g = (L - T)/k .
\end{array}
\]
\end{definition}

A single child under $\SBetween$ is centred, which is what the reference
implementation's comment says it intends.

\begin{warning}[\texttt{spacedBy} and \texttt{SPACE\_*} do not compose]
\label{sec:spacedquirk}
The gaps for the three $\mathsf{Space}$ arrangements are computed from $T$,
which \emph{excludes} \jv{spacedBy}, but the placement loop then still executes
\jv{tx += spacedBy}. A row or column that mixes \jv{spacedBy} with a
$\mathsf{Space}$ arrangement therefore overflows its container by exactly
$s\,(k-1)$. This is reproduced faithfully; see Appendix~\ref{app:examples},
Example~4.
\end{warning}

%%% ==================================================================
\section{Properties}
\label{sec:properties}

The Lean development proves the following. Unless noted, they are proved by
structural induction or by \jv{rfl} and are checked by the Lean kernel. They
live in \file{RemoteComposeLayout/Theorems.lean}.

\subsection{Placement preserves measurement}

Placement assigns offsets and never resizes.

\begin{property}[\texttt{rowPlace\_length}, \texttt{colPlace\_length}]
Placement returns exactly as many children as it was given.
\end{property}

\begin{property}[\texttt{rowPlace\_preserves\_widths}]
The list of widths is unchanged by row placement; dually
\texttt{colPlace\_preserves\_heights} for heights, and
\texttt{boxPlace\_preserves\_sizes} for both axes of a box.
\end{property}

\subsection{Main-axis progression}

\begin{property}[\texttt{rowPlace\_start\_head}]
Under $\Start$ with $s = 0$, the first child sits at $x = 0$.
\end{property}

\begin{property}[\texttt{rowPlace\_start\_tail}]
Under $\Start$, the remaining children are placed exactly as if the row had
begun at $\tau_1.w + s$.
\end{property}

\subsection{Cross-axis alignment}

\begin{property}[\texttt{crossX\_Start}, \texttt{crossY\_Top}]
$X(\Start, L, c) = 0$ and $Y(\Top, L, c) = 0$, for all $L$ and $c$.
\end{property}

\begin{property}[\texttt{crossX\_Center\_symmetric}]
$X(\Center, L, c) = \bigl(L - c\bigr)/2$ leaves equal margins:
the leading margin equals $L - c - X(\Center, L, c)$.
\end{property}

\begin{property}[\texttt{crossX\_End\_flush}, \texttt{crossY\_Bottom\_flush}]
Under $\End$ the child's trailing edge coincides with the container's:
$X(\End, L, c) + c = L$.
\end{property}

\begin{property}[\texttt{cross\_Center\_exact}]
A child exactly as large as its container is placed at $0$ under every
alignment.
\end{property}

\subsection{Arrangements fill their container}

\begin{property}[\texttt{spaceEvenly\_fills}]
Under $\SEvenly$ with $s = 0$, the $k$ children and the $k+1$ gaps exactly
partition $L$: $t_0 + T + k\,g = L$.
\end{property}

\begin{property}[\texttt{spaceAround\_fills}]
Under $\SAround$ with $s = 0$ and $k \neq 0$, the half-gaps at the ends and the
full gaps between children exactly partition $L$.
\end{property}

\begin{property}[\texttt{spaceBetween\_singleton}]
\[
  \Aa(\SBetween, L, C, T, 1) = \Aa(\Center, L, C, T, 1).
\]
\end{property}

\subsection{The modifier chain}

\begin{property}[\texttt{padding\_order\_irrelevant}]
$\Pp$ is invariant under permutation of the chain: accumulated padding does not
depend on where the padding modifiers sit.
\end{property}

\begin{property}[\texttt{definedWidth\_order\_matters}]
$\Dw$ \emph{is} order-dependent. Concretely,
\[
  \Dw\bigl(\pad(8)\cdot\wdt(\Exact(50))\bigr) = 66
  \;\neq\; 50 = \Dw\bigl(\wdt(\Exact(50))\cdot\pad(8)\bigr).
\]
This is a proved disequality, not an assumption: it pins the quirk of
Section~\ref{sec:orderquirk} down as a theorem.
\end{property}

\begin{property}[\texttt{definedWidth\_ignores\_height}]
A $\hgt$ modifier does not terminate the horizontal walk.
\end{property}

\begin{property}[\texttt{widthDim\_first\_wins}]
$\Dimw(\wdt(d,\iota)\cdot\wdt(d',\iota')\cdot\mu) = d$.
\end{property}

\subsection{Constraint enforcement}

\begin{property}[\texttt{clamp\_le}]
$\min(\max(x, \ell), u) \leq u$ for all $x, \ell, u$.
\end{property}

\begin{property}[\texttt{measure\_width\_le\_maxW}, \texttt{measure\_height\_le\_maxH}]
Under the hypotheses that the node declares no \jv{widthIn} (resp.\
\jv{heightIn}) and that its padding is non-negative, the measured size never
exceeds the incoming maximum:
$\Mm(n,\kappa).w \leq \kappa.w^+$.
\end{property}

\begin{remark}
The \jv{widthIn} hypothesis is necessary, not an artefact of the proof. Step~8
applies the \jv{widthIn} clamp \emph{after} the final constraint clamp, so a
\jv{widthIn} whose minimum exceeds the incoming maximum does push the measured
size past $\kappa.w^+$. The reference implementation has the same behaviour.
\end{remark}

\subsection{Collapsing and segmentation}

The two structural passes the collapsible and flow managers add are, on the face
of it, the two places in the algorithm where a child could go missing. These
properties say that neither loses one, and pin down the one respect in which
both are irreversible.

\begin{property}[\texttt{collapseSweep\_length}]
The collapsing sweep reaches a decision about every child, exactly once:
$\len\bigl(\Sigma(B, \vec{x}, \alpha, f)\bigr) = \len(\vec{x})$.
\end{property}

\begin{property}[\texttt{collapseSweep\_overflow}]
\emph{Refusal is sticky.} Once the sweep has refused one child, it refuses every
child after it in priority order, however small:
\[
  \Sigma(B, \vec{x}, \alpha, \textsf{true})
    = \bigl[\,(i, \textsf{true}) \;\big|\; (i, \tau) \in \vec{x}\,\bigr].
\]
A collapsible manager never back-fills.
\end{property}

\begin{property}[\texttt{applyFlags\_preserves\_rects}]
Applying the decisions changes visibility and nothing else: the list of
rectangles is unaffected. Collapsing does not move or resize anything; the
consequences of the decision are drawn later, by the placement sweep.
\end{property}

\begin{property}[\texttt{priorityLt\_irrefl}]
$\prec_o$ is irreflexive --- no child precedes itself. This is the least a sort
may ask of a comparator, and it is not automatic here, since $\prec_o$ is
defined by the truncating \jv{(int) (priority2 - priority1)} rather than by a
plain comparison.
\end{property}

\begin{property}[\texttt{priorityInsert\_length}, \texttt{sortWithPriorities\_length}]
The priority sort keeps every child: $\len(\mathrm{sort}_o(\vec{n})) =
\len(\vec{n})$. It decides only the order in which children are offered space;
which of them survive is the sweep's business.
\end{property}

\begin{property}[\texttt{flowSegment\_length}]
Segmentation assigns every child either a row or a refusal, and never both or
neither: $\len(\mathrm{seg}(\ldots, \vec{e})) = \len(\vec{e})$.
\end{property}

\begin{property}[\texttt{flowRowCount\_pos}]
A flow always opens at least one row, even when every child is refused and even
when there are none. The count starts at $1$ because the reference
implementation appends an empty row before it begins walking --- which is
exactly why $k_{\mathrm{lines}}$ bites a row earlier than one would expect.
\end{property}

\begin{property}[\texttt{maskRow\_length}]
Masking a row hides the other rows and disturbs nothing else, preserving the
length of the child list. This is what makes the masking construction of
Section~\ref{sec:flow} faithful to the reference implementation's habit of
passing one row's sublist to \jv{RowLayout}.
\end{property}

%%% ==================================================================
\section{Layout Optimization Passes and Invariants}
\label{sec:optimizations}

The reference Java player defines four core layout optimizations in
\jv{CoreDocument.java}, governed by a four-bit mask:
\begin{align*}
  \jv{OPTIMIZATION\_NONE}               &= 0, \\
  \jv{OPTIMIZATION\_MEASURE\_CACHE}       &= 1 \quad (2^0), \\
  \jv{OPTIMIZATION\_LAYOUT\_BOUNDARIES}  &= 2 \quad (2^1), \\
  \jv{OPTIMIZATION\_FLAT\_MEASURE\_PASS} &= 4 \quad (2^2), \\
  \jv{OPTIMIZATION\_CONSTRAINTS\_CACHE}  &= 8 \quad (2^3), \\
  \jv{OPTIMIZATION\_ALL}                 &= 15.
\end{align*}
The Lean module \file{RemoteComposeLayout/Optimizations.lean} models each of
these passes, their interaction with document invalidation and subtree
remeasurement, and proves their correctness invariants.

\subsection{Optimization levels and bitmask isomorphism}

The flags are represented as a four-field record \lo{OptimizationLevel}:
\[
  \lo{OptimizationLevel} = \langle \textsf{measureCache}, \textsf{layoutBoundaries},
  \textsf{flatMeasurePass}, \textsf{constraintsCache} \rangle.
\]
The bitmask packing and unpacking functions satisfy an exact roundtrip isomorphism:

\begin{property}[\lo{bitmask\_roundtrip}]
For any configuration $\mathrm{opt} \in \lo{OptimizationLevel}$,
\[
  \mathrm{fromNat}(\mathrm{toNat}(\mathrm{opt})) = \mathrm{opt}.
\]
Every combination of optimization flags is uniquely and reversibly encoded in the
range $[0, 15]$.
\end{property}

\subsection{Relayout boundaries and incremental layout}
\label{sec:opt-boundaries}

A component is declared a \emph{relayout boundary}
(\jv{Component.isRelayoutBoundary}) when its dimensions are externally
determined and cannot depend on the size of its children:
\begin{definition}[Relayout Boundary]
A node $n$ is a relayout boundary if and only if:
\[
  \mathrm{isBoundary}(n) \iff
  \bigl(\Dimw(n).\mathrm{isExact} \land \Dimh(n).\mathrm{isExact}\bigr) \;\lor\;
  \bigl(\Dimw(n).\mathrm{isFill} \land \Dimh(n).\mathrm{isFill}\bigr).
\]
\end{definition}

When a descendant component triggers an invalidation via \jv{invalidateMeasure()},
the tree climbs parent pointers:
\begin{enumerate}
  \item If \jv{OPTIMIZATION\_LAYOUT\_BOUNDARIES} is disabled, invalidation ascends to
  the document root, forcing a full pass (\lo{InvalidationTarget.fullRoot}).
  \item If enabled, the ascent halts at the first ancestor satisfying
  $\mathrm{isBoundary}(n)$, registering that ancestor's path as a dirty boundary
  (\lo{InvalidationTarget.dirtyBoundary}).
\end{enumerate}

During playback, \jv{RootLayoutComponent.performPartialLayoutPass} traverses only
registered dirty boundaries:
\[
  \tau' = \mathrm{layout}(n_{\mathrm{dirty}}, \tau.\mathrm{w}, \tau.\mathrm{h})
  \quad\text{with}\quad
  (x', y') = (\tau.x, \tau.y).
\]
The surrounding layout tree, sibling positions, and ancestor geometry are
preserved without re-execution.

\begin{property}[\lo{exact\_boundary\_width\_fixed}]
An exact-sized component's boundary status is an invariant under arbitrary child
mutations: for any child lists $\vec{c}_1, \vec{c}_2$, letting
\begin{gather*}
  n_1 = \lo{Node.box}([\wdt(\Exact(w)), \hgt(\Exact(h))], \vec{c}_1), \\
  n_2 = \lo{Node.box}([\wdt(\Exact(w)), \hgt(\Exact(h))], \vec{c}_2),
\end{gather*}
we have $\mathrm{isBoundary}(n_1) = \mathrm{isBoundary}(n_2) = \textsf{true}$.
Mutations within the subtree cannot disturb the boundary classification.
\end{property}

\subsection{Constraint caching and memoization}
\label{sec:opt-cache}

\jv{ComponentMeasure.hasCachedConstraints} memoizes measurement results. A cache
hit occurs if the incoming constraints match one of three predicates:
\begin{definition}[Constraint Cache Hit Criteria]
Let $C = (w_{\min}, w_{\max}, h_{\min}, h_{\max})$ be the cached constraints,
$(w^*, h^*)$ the previous dimensions, and $C' = (w'_{\min}, w'_{\max},$
$h'_{\min}, h'_{\max})$ the incoming constraints. Then $C'$ hits the cache if:
\begin{align*}
  \mathrm{Match}_1 &\iff C' = C \quad \text{(exact match)}, \\
  \mathrm{Match}_2 &\iff w'_{\min} = w_{\min} \;\land\; w'_{\max} = w_{\max} \;\land\;
  h'_{\min} = h^* = h'_{\max}, \\
  \mathrm{Match}_3 &\iff h'_{\min} = h_{\min} \;\land\; h'_{\max} = h_{\max} \;\land\;
  w'_{\min} = w^* = w'_{\max}.
\end{align*}
\end{definition}

\begin{property}[\lo{cache\_exact\_match}]
If incoming constraints exactly match the cached descriptor,
$\lo{hasCachedConstraints}$ evaluates unconditionally to $\textsf{true}$.
\end{property}

\subsection{Generational flat-array measure pass}
\label{sec:opt-flatpass}

Rather than allocating per-component \jv{ComponentMeasure} objects and clearing
them with an $O(N)$ tree walk, \jv{FlatMeasurePass.java} assigns each component an
integer index into preallocated flat arrays. Invalidation is achieved in $O(1)$
via a monotonically increasing generation token:
\[
  g \gets g + 1.
\]
A slot $i$ is fresh for the current layout cycle if and only if:
\[
  \mathrm{isFresh}(i) \iff \mathrm{slot}[i].\mathrm{generation} = g.
\]
When $g$ reaches $2^{31}-1$, the generation resets to $1$ and the array is cleared.

\begin{property}[\texttt{flatPass\_clear\_expires\_slots}]
For any state where all allocated slots have generation $\leq g < 2^{31}-1$,
incrementing the pass generation guarantees that every existing slot becomes stale
immediately:
\[
  \mathrm{isFresh}(\lo{clear}(s), i) = \textsf{false} \quad (\forall i).
\]
No slot memory needs to be zeroed or reallocated during steady-state frame rendering.
\end{property}

\subsection{Flex two-pass short-circuiting}
\label{sec:opt-flex}

In \jv{RowLayout} and \jv{ColumnLayout}, allocating extra space to weighted
children requires a second measurement sweep. The engine short-circuits this
when no weighted children are present:

\begin{property}[\texttt{unweighted\_row\_single\_pass}]
If no child declares a weight on either dimension:
\[
  \bigvee_{c \in \vec{c}} \bigl(\Dimw(c).\mathrm{hasWeight} \lor \Dimh(c).\mathrm{hasWeight}\bigr) = \textsf{false},
\]
then all children satisfy:
\[
  \bigwedge_{c \in \vec{c}} \bigl(\neg \Dimw(c).\mathrm{hasWeight} \land \neg \Dimh(c).\mathrm{hasWeight}\bigr) = \textsf{true}.
\]
The measure pass finishes in a single pass without invoking flex space redistribution.
\end{property}

\subsection{Measure policy hierarchy}
\label{sec:opt-policies}

Layout behavior across RemoteCompose versions is formalized via five progressive
policy configurations:
\[
  \lo{MeasurePolicyVersion} = \left\{
  \begin{aligned}
    &\lo{v0\_legacy}, \lo{v1\_baseModern}, \lo{v2\_insetWrap}, \\
    &\lo{v3\_inlineExpression}, \lo{v4\_enforceConstraints}
  \end{aligned}
  \right\}.
\]
Each configuration determines three behavioral predicates:
\[
  \lo{MeasurePolicyConfig} = \langle \textsf{applyInsetWrap},
  \textsf{updateComponentValues}, \textsf{enforceConstraints} \rangle.
\]
The policies form a monotonic progression:
\begin{align*}
  \lo{v0\_legacy}             &\mapsto \langle \textsf{false}, \textsf{false}, \textsf{false} \rangle, \\
  \lo{v1\_baseModern}         &\mapsto \langle \textsf{false}, \textsf{false}, \textsf{false} \rangle, \\
  \lo{v2\_insetWrap}          &\mapsto \langle \textsf{true},  \textsf{false}, \textsf{false} \rangle, \\
  \lo{v3\_inlineExpression}   &\mapsto \langle \textsf{true},  \textsf{true},  \textsf{false} \rangle, \\
  \lo{v4\_enforceConstraints} &\mapsto \langle \textsf{true},  \textsf{true},  \textsf{true}  \rangle.
\end{align*}

%%% ==================================================================
\section{Extended Modifiers and Transformations}
\label{sec:extended-modifiers}

The Lean module \file{RemoteComposeLayout/ExtendedModifiers.lean} formalizes
the full suite of geometric transformations, visibility states, alignment anchors,
and rendering properties defined across \file{remote-core/operations/layout/modifiers/}.

\subsection{Visibility modes and space preservation}
\label{sec:mod-visibility}

The reference player defines component visibility via the operation
\jv{ComponentVisibilityOperation} (\jv{Component.java}), taking values in:
\[
  \lo{VisibilityMode} = \{\lo{visible},\, \lo{invisible},\, \lo{gone}\}.
\]
These modes govern two distinct semantic questions: whether the component occupies
spatial layout bounds, and whether it produces visual raster output or receives touch:
\begin{align*}
  \lo{takesSpace}(v) &= \begin{cases}
    \textsf{true}  & \text{if } v \in \{\lo{visible}, \lo{invisible}\}, \\
    \textsf{false} & \text{if } v = \lo{gone},
  \end{cases} \\
  \lo{isPainted}(v) &= \begin{cases}
    \textsf{true}  & \text{if } v = \lo{visible}, \\
    \textsf{false} & \text{if } v \in \{\lo{invisible}, \lo{gone}\}.
  \end{cases}
\end{align*}
The effective layout dimension $\lo{effectiveDim}(v, s)$ of a component with
natural dimension $s \in \Sc$ is defined by:
\[
  \lo{effectiveDim}(v, s) = \begin{cases}
    s   & \text{if } \lo{takesSpace}(v), \\
    0.0 & \text{otherwise.}
  \end{cases}
\]

\begin{property}[\lo{invisible\_preserves\_space}]
An invisible component allocates identical spatial dimensions to a fully visible component:
\[
  \lo{effectiveDim}(\lo{invisible}, s) = \lo{effectiveDim}(\lo{visible}, s) = s \quad (\forall s \in \Sc).
\]
Invisible items maintain line breaks, cross-axis height, and flex spacing without drawing.
\end{property}

\begin{property}[\lo{gone\_collapses\_space}]
A gone component unconditionally collapses its spatial footprint to zero:
\[
  \lo{effectiveDim}(\lo{gone}, s) = 0.0 \quad (\forall s \in \Sc).
\]
It is treated as non-existent during measure and placement loops.
\end{property}

\subsection{Typographic baseline alignment}
\label{sec:mod-baseline}

In horizontal layouts, text and icon components can be aligned relative to their
typographic baseline via \jv{AlignByModifierOperation}. A baseline descriptor
$\lo{BaselineDesc} = \langle h, b \rangle$ records the child height $h$ and the
offset $b \in \Sc$ from the top of the component to the text baseline.

When a row container computes cross-axis placement for baseline-aligned children,
it determines the maximal baseline across all aligned siblings:
\[
  B_{\max} = \max_{i} b_i.
\]
Each child is then assigned a cross-axis placement offset:
\[
  y_i = \lo{alignByOffsetY}(B_{\max}, b_i) \triangleq B_{\max} - b_i.
\]

\begin{property}[\lo{baseline\_alignment\_coincident}]
For any child with baseline $b_i$, its global baseline coordinate in the container satisfies:
\[
  y_i + b_i = (B_{\max} - b_i) + b_i = B_{\max}.
\]
All aligned children achieve exact typographic baseline coincidence regardless of differences
in font size, ascent, or total component height.
\end{property}

\subsection{Visual translation offset}
\label{sec:mod-offset}

The \jv{OffsetModifierOperation} applies a local displacement vector
$\langle dx, dy \rangle \in \Sc^2$ to the placed bounding box:
\[
  \lo{applyToRect}(\langle dx, dy \rangle, \langle x, y, w, h \rangle) \triangleq
  \langle x + dx, y + dy, w, h \rangle.
\]

\begin{property}[\lo{offset\_preserves\_dimensions}]
Local offset translation preserves rectangle dimensions invariant:
\[
  (\lo{applyToRect}(\mathrm{off}, r)).w = r.w \;\land\;
  (\lo{applyToRect}(\mathrm{off}, r)).h = r.h \quad (\forall \mathrm{off}, r).
\]
Local offsets shift visual rendering and semantic hit-testing without perturbing
sibling layout coordinates or container boundary measurements.
\end{property}

\subsection{Scrollable viewports and bounds clamping}
\label{sec:mod-scroll}

\jv{ScrollModifierOperation} manages virtual scrollable viewports over child content
larger than the allocated window bounds. For scroll position $p$ and maximum scrollable
extent $M \in \Sc$, the player clamps the translation:
\[
  \lo{clampScroll}(p, M) \triangleq \begin{cases}
    0.0 & \text{if } M \le 0.0 \lor p \le 0.0, \\
    M   & \text{if } p \ge M, \\
    p   & \text{otherwise.}
  \end{cases}
\]

\begin{property}[\lo{scroll\_clamp\_nonneg} and \lo{scroll\_clamp\_upper}]
For any query position $p$ and non-negative extent $M \ge 0.0$:
\[
  0.0 \le \lo{clampScroll}(p, M) \le M.
\]
Virtual scrolling cannot expose unpainted negative coordinates or exceed valid content extents.
\end{property}

\subsection{Graphics layers and Z-index ordering}
\label{sec:mod-graphics}

Visual styling and affine transformations are encapsulated in \jv{GraphicsLayerModifierOperation},
represented by the record $\lo{GraphicsLayer}$:
\[
  \lo{GraphicsLayer} = \langle s_x, s_y, \theta, t_x, t_y, \alpha \rangle
\]
with default identity values $s_x = s_y = \alpha = 1.0$ and $\theta = t_x = t_y = 0.0$.
\begin{property}[\lo{identity\_is\_identity}]
The default $\lo{GraphicsLayer.identity}$ satisfies $\lo{isIdentity}$ unconditionally.
\end{property}

Component stacking order across overlapping siblings is parameterized by \jv{ZIndexModifierOperation},
recording a scalar $z \in \Sc$. Components with higher $z$-values are painted in front and receive
interactive hit-test events ahead of lower-$z$ siblings.

%%% ==================================================================
\section{Extended Layout Containers}
\label{sec:extended-containers}

The Lean module \file{RemoteComposeLayout/ExtendedContainers.lean} formalizes
three specialized layout managers implemented in \file{remote-core/operations/layout/managers/}:
\jv{FitBoxLayout}, \jv{StateLayout}, and \jv{ImageLayout}.

\subsection{FitBox layout: greedy alternative selection}
\label{sec:fitbox}

\jv{FitBoxLayout} acts as an adaptive selector among a list of alternative candidate
representations $[c_0, c_1, \ldots, c_{m-1}]$:
\begin{enumerate}
  \item Candidates are evaluated sequentially in declaration order.
  \item The \emph{first} candidate $c_k$ whose measured dimensions satisfy
    $c_k.w \le w_{\max}$ and $c_k.h \le h_{\max}$ is admitted as the visible child.
  \item All other candidates are collapsed to $\lo{gone}$.
\end{enumerate}

\begin{property}[\lo{fitbox\_at\_most\_one\_visible}]
For any constraints $(w_{\max}, h_{\max})$ and candidate list $\vec{c}$:
\[
  \bigl|\{r \in \lo{fitBoxSelect}(w_{\max}, h_{\max}, \vec{c}) \mid r.\mathrm{visible} = \textsf{true}\}\bigr| \le 1.
\]
FitBox guarantees mutual exclusion among candidate representations.
\end{property}

\begin{property}[\lo{fitbox\_admitted\_fits\_container}]
\leavevmode\\
If a candidate $r$ is admitted with $r.\mathrm{visible} = \textsf{true}$, it strictly fits
within the container constraints:
\[
  r.w \le w_{\max} \;\land\; r.h \le h_{\max}.
\]
\end{property}

\subsection{State layout: declarative state switching}
\label{sec:statelayout}

\jv{StateLayout} models a dynamic multi-state container indexed by an integer variable
$k \in \mathbb{N}$ over an array of subtrees $[s_0, s_1, \ldots, s_{m-1}]$:
\[
  \lo{resolveStateLayout}(\langle k, \vec{s} \rangle) \triangleq \vec{s}[k].
\]
Unselected states are omitted from the composition tree.

\begin{property}[\lo{state\_isolation}]
If two configurations $\vec{s}_1, \vec{s}_2$ share active state $k$ with identical subtree $t$:
\[
  \vec{s}_1[k] = \vec{s}_2[k] = \mathrm{some}(t) \implies
  \lo{resolveStateLayout}(\langle k, \vec{s}_1 \rangle) = \lo{resolveStateLayout}(\langle k, \vec{s}_2 \rangle).
\]
Mutations to inactive states cannot alter the current visual layout.
\end{property}

\subsection{Image layout and aspect ratio scaling}
\label{sec:imagelayout}

\jv{ImageLayout} sizes an image with intrinsic content dimensions $(w_i, h_i)$ inside container
bounds $(w_b, h_b)$ under three scaling policies:
\begin{itemize}
  \item \lo{fillBounds}: Direct stretch to $(w_b, h_b)$, ignoring aspect ratio.
  \item \lo{fit}: Uniform scaling by $s = \min(w_b / w_i, h_b / h_i)$ yielding $(w_i \cdot s, h_i \cdot s)$.
  \item \lo{inside}: If $(w_i, h_i)$ fits inside $(w_b, h_b)$, retains $(w_i, h_i)$; otherwise scales as \lo{fit}.
\end{itemize}

\begin{property}[\lo{fit\_preserves\_aspect\_ratio}]
Under $\lo{fit}$ scaling with $s \ne 0$ and $h_i \ne 0$:
\[
  \frac{w_i \cdot s}{h_i \cdot s} = \frac{w_i}{h_i}.
\]
The intrinsic geometric aspect ratio is preserved under uniform scaling.
\end{property}

%%% ==================================================================
\section{Window Hierarchy and Pointer Hit Testing}
\label{sec:window-hit}

The Lean module \file{RemoteComposeLayout/WindowAndHitTest.lean} formalizes
the transformation from local component coordinates to global window space
and the dispatch of interactive pointer events.

\subsection{Window coordinate accumulation}
\label{sec:window-loc}

Each component's placement $(x, y)$ is defined relative to its direct parent.
The reference method \jv{Component.getLocationInWindow} ascends the ancestor chain
$\vec{f} = [f_1, \ldots, f_k]$, where each frame records local origin $(f.x, f.y)$
and scroll offsets $(f.\mathrm{scrollX}, f.\mathrm{scrollY})$:
\[
  \lo{accumulateWindowLocation}(\vec{f}, x, y) \triangleq
  \biggl(x + \sum_{f \in \vec{f}} (f.x - f.\mathrm{scrollX}),\;
         y + \sum_{f \in \vec{f}} (f.y - f.\mathrm{scrollY})\biggr).
\]

\begin{property}[\lo{subtree\_locality\_shift}]
Shifting each ancestor frame origin by $(\Delta x, \Delta y)$ shifts the resulting
global window location of all descendant elements by exactly:
\[
  \lo{accumulateWindowLocation}(\vec{f}', x, y) =
  \lo{accumulateWindowLocation}(\vec{f}, x, y) + |\vec{f}| \cdot (\Delta x, \Delta y).
\]
Hierarchical translation is strictly linear in the depth of the ancestor chain.
\end{property}

\subsection{Soundness of pointer hit-testing}
\label{sec:hittest}

Pointer events $(px, py)$ are dispatched by \jv{TouchHandler} and \jv{ClickHandler}
down the tree of interactive targets $\lo{HitTarget} = \lo{node}(id, r, v, z, \vec{c})$:
\begin{enumerate}
  \item A node is rejected if $\neg\lo{takesSpace}(v)$ or $\neg\lo{pointInRect}(px, py, r)$, where:
    \[
      \lo{pointInRect}(px, py, r) \iff r.x \le px < r.x + r.w \;\land\; r.y \le py < r.y + r.h.
    \]
  \item If the point is geometrically contained, children are tested recursively
    in front-to-back order (deepest child takes priority).
  \item If no child accepts the event and the node satisfies $\lo{isPainted}(v)$,
    the event is captured by the parent node $id$.
\end{enumerate}

\begin{property}[\lo{hitTest\_soundness}]
If $\lo{hitTest}(T, px, py) = \mathrm{some}(T.id)$, the query point $(px, py)$
is guaranteed to lie within the bounding box of target $T$:
\[
  T.r.x \le px < T.r.x + T.r.w \;\land\; T.r.y \le py < T.r.y + T.r.h.
\]
A component can never capture pointer events occurring outside its spatial layout bounds.
\end{property}

%%% ==================================================================
\section{Layout Animations and Continuity}
\label{sec:animation}

The Lean module \file{RemoteComposeLayout/Animation.lean} formalizes smooth
temporal layout transitions defined in \jv{AnimateMeasure.java}.

\subsection{Temporal progress mapping}
\label{sec:anim-progress}

For animation elapsed time $t$ and total duration $T \in \Sc$, the normalized
progress $\alpha(t, T) = \lo{animProgress}(t, T)$ is computed by:
\[
  \lo{animProgress}(t, T) \triangleq \begin{cases}
    1.0   & \text{if } T \le 0.0, \\
    0.0   & \text{if } t \le 0.0, \\
    1.0   & \text{if } t \ge T, \\
    t / T & \text{otherwise.}
  \end{cases}
\]

\begin{property}[\lo{animProgress\_nonneg} and \lo{animProgress\_le\_one}]
For all $t, T \in \Sc$, the normalized animation progress is bounded:
\[
  0.0 \le \lo{animProgress}(t, T) \le 1.0.
\]
Progress cannot underflow before the start instant or overflow after completion.
\end{property}

\subsection{Geometric bounds interpolation and continuity}
\label{sec:anim-interp}

During layout transitions, component rectangles are interpolated continuously between
original bounds $R_0 = \langle x_0, y_0, w_0, h_0 \rangle$ and target bounds
$R_1 = \langle x_1, y_1, w_1, h_1 \rangle$ using scalar linear interpolation
$\lo{lerp}(a, b, p) = a + p(b - a)$:
\[
  \lo{interpRect}(R_0, R_1, p) \triangleq \langle
    \lo{lerp}(x_0, x_1, p),\, \lo{lerp}(y_0, y_1, p),\,
    \lo{lerp}(w_0, w_1, p),\, \lo{lerp}(h_0, h_1, p) \rangle.
\]

\begin{property}[\lo{animate\_at\_zero}]
At $t = 0$ with duration $T > 0$, the interpolated rectangle matches the origin bounds exactly:
\[
  \lo{interpRect}(R_0, R_1, \lo{animProgress}(0, T)) = R_0.
\]
Transitions begin with zero discontinuous jumping.
\end{property}

\begin{property}[\lo{animate\_at\_end}]
At any time $t \ge T$ with $T > 0$, the interpolated rectangle matches the target bounds exactly:
\[
  \lo{interpRect}(R_0, R_1, \lo{animProgress}(t, T)) = R_1.
\]
Transitions settle precisely on their target layout without residual drift.
\end{property}

\begin{property}[\lo{animate\_stationary}]
If origin and target bounds coincide ($R_0 = R_1 = R$), the rectangle is stationary:
\[
  \lo{interpRect}(R, R, p) = R \quad (\forall p \in \Sc).
\]
\end{property}

%%% ==================================================================
\section{Dynamic Expressions and Runtime Invariants}
\label{sec:dynamic-expressions}

The Lean module \file{RemoteComposeLayout/DynamicExpressions.lean} formalizes
dynamic layout expressions evaluated at runtime from state variables:
\begin{center}
\file{remote-core/operations/expressions/}.
\end{center}

\subsection{Expression syntax and evaluation}
\label{sec:expr-eval}

Layout expressions $\lo{LayoutExpr}$ are formed according to the grammar:
\[
  e ::= \lo{const}(v) \mid \lo{var}(i) \mid e_1 + e_2 \mid e_1 - e_2 \mid e_1 \cdot e_2 \mid \lo{clamp}(e, \min, \max)
\]
where $v, \min, \max \in \Sc$ and $i \in \mathbb{N}$ denotes a state variable ID.
Given an evaluation environment $\eta : \mathbb{N} \to \Sc$, the evaluation function
$\lo{evalExpr}(\eta, e)$ computes the scalar value, interpreting $\lo{clamp}(e, v_{\min}, v_{\max})$ as:
\[
  \lo{evalExpr}(\eta, \lo{clamp}(e, v_{\min}, v_{\max})) \triangleq \max\bigl(v_{\min}, \min(\lo{evalExpr}(\eta, e), v_{\max})\bigr).
\]

\subsection{Safety and boundedness invariants}
\label{sec:expr-invariants}

Dynamic layout dimensions are subject to verified safety bounds:

\begin{property}[\lo{eval\_const}]
Constant expressions evaluate independently of the variable environment:
\[
  \lo{evalExpr}(\eta_1, \lo{const}(v)) = \lo{evalExpr}(\eta_2, \lo{const}(v)) = v \quad (\forall \eta_1, \eta_2).
\]
\end{property}

\begin{property}[\lo{eval\_clamp\_lower} and \lo{eval\_clamp\_upper}]
For any expression $e$, environment $\eta$, and bounds $v_{\min} \le v_{\max}$:
\[
  v_{\min} \le \lo{evalExpr}(\eta, \lo{clamp}(e, v_{\min}, v_{\max})) \le v_{\max}.
\]
Dynamic dimensions constrained by clamping cannot violate container min/max boundaries.
\end{property}

\begin{property}[\lo{eval\_env\_agree}]
Let $\mathrm{vars}(e) \subset \mathbb{N}$ be the set of free variables occurring in $e$.
If two environments $\eta_1, \eta_2$ agree on all free variables of $e$:
\[
  \bigl(\forall i \in \mathrm{vars}(e),\; \eta_1(i) = \eta_2(i)\bigr) \implies
  \lo{evalExpr}(\eta_1, e) = \lo{evalExpr}(\eta_2, e).
\]
Evaluation depends only on referenced variables; unrelated state updates cannot perturb layout expressions.
\end{property}

%%% ==================================================================
\section{Autosize Text and Repaint Loop Invariants}
\label{sec:autosize-loop}

The Lean module \file{RemoteComposeLayout/AutosizeRepaintLoop.lean} formalizes
the interplay between adaptive text sizing and collapsible layout containers,
and mathematically proves the root cause of the infinite repaint loop bug reported
in production.

\subsection{Autosize text sizing dynamics}
\label{sec:autosize-dynamics}

A text node $\ttext(\mu, \theta)$ with autosizing enabled (\jv{CoreText.java}, lines 705--757)
determines its font size $s \in [s_{\min}, s_{\max}]$ dynamically by finding the maximal
font size whose rendered bounding box fits within the incoming container constraints.
Formally, let $h(s)$ be the monotonically non-decreasing height of the text at font size $s$:
\[
  s_1 \le s_2 \implies h(s_1) \le h(s_2).
\]
Given an upper height constraint $H^+$, the autosizing font selection function is:
\[
  \lo{selectAutosizeFont}(\theta, H^+) \triangleq \begin{cases}
    s_{\max} & \text{if } H^+ \ge h(s_{\max}), \\
    s_{\min} & \text{otherwise.}
  \end{cases}
\]

\begin{property}[\lo{autosize\_balloons\_under\_fltMax}]
When an autosizing text component is measured against an unbounded height budget $H^+ = \Om = \lo{fltMax}$:
\[
  \lo{selectAutosizeFont}(\theta, \Om) = s_{\max} \quad\text{and}\quad
  \lo{autosizeMeasureHeight}(\theta, \Om) = h(s_{\max}).
\]
Because $\Om$ is unbinding, the component unconditionally balloons to its maximum font size.
\end{property}

\subsection{False overflow and container self-collapse}
\label{sec:autosize-collapse}

In \jv{CollapsibleColumnLayout}, measurement proceeds in two phases (Section~\ref{sec:ccol}):
\begin{enumerate}
  \item \textbf{Phase 1 (\jv{computeVisibleChildren})}: Each child is provisionally measured.
    Crucially, in the reference implementation, unweighted children are queried with
    $H^+ = \Om$ (\jv{Float.MAX\_VALUE}).
  \item \textbf{Phase 2 (\jv{collapseSweep})}: The measured heights from Phase 1 are tested
    against the actual container budget $B$.
\end{enumerate}

\begin{property}[\lo{collapsible\_child\_falsely\_collapsed}]
Suppose an autosizing text component fits comfortably at minimum size ($h(s_{\min}) \le B$),
but exceeds the container budget at maximum size ($h(s_{\max}) > B$).
Because Phase~1 measured the child with $H^+ = \Om$, Phase~2 receives the ballooned height
$h(s_{\max}) > B$ and marks the child $\gone = \textsf{true}$.
\end{property}

\begin{property}[\lo{collapsible\_column\_total\_collapse}]
If all children are marked $\gone$ by the collapsing sweep:
\[
  \lo{countVisible}(\lo{applyFlags}(\vec{\tau}, \vec{\gone})) = 0.
\]
The collapsible column marks itself $\gone$ and reports a measured height of $0$.
\end{property}

\subsection{Infinite repaint loop as a dynamical system}
\label{sec:autosize-repaint-loop}

The interaction between the layout engine and the Android view lifecycle
(\jv{RemoteComposeView.java} and \jv{CoreDocument.java}) is formalized as a discrete
dynamical system on the frame state $\sigma = \langle V_h, R_h, m, p, k \rangle$, where
$V_h$ is the host viewport height, $R_h$ is the root component's measured height,
$m \in \{\textsf{true}, \textsf{false}\}$ indicates $\lo{needsMeasure}$,
$p \in \{\textsf{true}, \textsf{false}\}$ indicates $\lo{needsRepaint}$, and
$k \in \mathbb{N}$ is the $\lo{repaintNext}$ frame counter.

A frame tick $\Phi(\sigma) = \lo{paintStep}(\lo{measureStep}(\sigma))$ consists of:
\begin{enumerate}
  \item \textbf{Measure step}: If $m = \textsf{true}$, the host re-measures the root.
    Under the autosize bug, the root collapses again to $R_h = 0$, setting $m \gets \textsf{false}$.
  \item \textbf{Paint step}: \jv{CoreDocument.paint} checks whether the viewport dimension matches the root:
    \[
      \lo{dirty} \triangleq (V_h \neq R_h).
    \]
    If $\lo{dirty}$, it invokes \jv{invalidateMeasure()}, which asserts $m \gets \textsf{true}$
    and $p \gets \textsf{true}$. When $p = \textsf{true}$, it posts an immediate frame callback by
    setting $k \gets 1$.
\end{enumerate}

\begin{property}[\lo{infinite\_repaint\_loop\_invariant}]
For any host viewport $V_h > 0$ with an initial collapsed root $R_h = 0$, the state under iteration
$\Phi^n(\sigma_0)$ satisfies for all $n \ge 1$:
\[
  \Phi^n(\sigma_0) = \langle V_h, 0, \textsf{true}, \textsf{true}, 1 \rangle,
  \qquad \lo{repaintNext} = 1 \;\land\; \lo{needsRepaint} = \textsf{true}.
\]
The non-quiescent state is an attractor with cycle period $1$. The system never quiesces,
triggering an infinite stream of 60/120~Hz Choreographer redraws.
\end{property}

\subsection{Sound fix and verified quiescence}
\label{sec:autosize-fix}

The Lean development proves two theorems establishing the formal fix:
\begin{enumerate}
  \item \textbf{Sound child budget (\lo{sound\_autosize\_fits})}: In Phase~1, the container must pass
    its remaining container budget $B_{\mathrm{rem}} \le B$ rather than $\Om$. Then:
    \[
      \lo{autosizeMeasureHeight}(\theta, B_{\mathrm{rem}}) \le B.
    \]
    The child fits within the container and is not falsely collapsed.
  \item \textbf{Quiescence verification (\lo{sound\_engine\_quiescence})}: Invalidation in \jv{paint}
    must check whether \emph{viewport constraints} changed, rather than comparing viewport size against
    content wrap height. When viewport constraints are stable, the system reaches the quiescent state
    $\langle V_h, R_h, \textsf{false}, \textsf{false}, 0 \rangle$ in exactly one frame.
\end{enumerate}

%%% ==================================================================
\section{Conformance}

\label{sec:conformance}

\subsection{Method}

The specification is validated against the project's existing gold corpus,
which lives at
\begin{center}
\file{compose/remote/specification/conformance/\{tests,gold\}/layout/}
\end{center}
and is generated from the Java player itself. A generator script,
\file{tools/gen\_conformance.py}, reads each test document and its gold file,
translates the component tree into the syntax of Section~3, and emits Lean
\jv{\#guard} assertions comparing $\mathrm{abs}(\mathrm{layout}(n, W, H))$
against the recorded rectangles. The assertions are discharged by the Lean
kernel when the library is built, so a regression breaks the build.

Two details of the gold format matter. Its \jv{expected\_tree} is stored in
\emph{reverse} pre-order, and its first element after reversal is a
\jv{RootLayoutComponent} viewport wrapper that must be dropped. Comparison uses
the corpus' own tolerance, $\varepsilon = \tfrac12$ pixel per coordinate.

Every check compares \emph{two} lists: the rectangles, against the gold
\jv{x}/\jv{y}/\jv{width}/\jv{height}, and the visibility flags $g$, against the
gold \jv{isGone}. A specification that placed every child correctly but
collapsed the wrong one would fail. The assertion form is \lo{conformsVis}.

Documents that use any construct outside the scope of Section~1.3 are skipped
rather than approximated, so that a passing run means what it says.

\subsection{Results}

\begin{center}
\begin{tabular}{@{}lr@{}}
\toprule
Documents in scope & \textbf{99} \\
Viewport configurations checked & \textbf{294} \\
Geometry failures & \textbf{0} \\
Visibility failures & \textbf{0} \\
Documents skipped as out of scope & 75 \\
\bottomrule
\end{tabular}
\end{center}

\noindent The skipped documents break down as follows.

\begin{center}
\begin{tabular}{@{}rl@{\qquad}rl@{}}
\toprule
\multicolumn{2}{c}{\textbf{Reason}} & \multicolumn{2}{c}{\textbf{Reason}} \\
\midrule
14 & \jv{stateLayout}             &  2 & modifier \jv{animationSpec} \\
10 & \jv{text}                    &  2 & \jv{widthIn} without a width modifier \\
10 & \jv{fitBox}                  &  2 & malformed component \\
 8 & modifier \jv{offset}         &  1 & \jv{canvas} \\
 8 & modifier \jv{visibility}     &  1 & \jv{image} \\
 7 & modifier \jv{verticalScroll} &  1 & \jv{heightIn} without a height modifier \\
 3 & modifier \jv{horizontalScroll} & 1 & modifier \jv{alignByBaseline} \\
 3 & column alignment             &  1 & modifier \jv{fillParentMaxWidth} \\
   &                              &  1 & modifier \jv{wrapContentSize} \\
\bottomrule
\end{tabular}
\end{center}

The three collapsible and flow entries that dominated this table in the
Box/Row/Column-only edition --- 18 \jv{collapsibleColumn}, 18
\jv{collapsibleRow}, 15 \jv{flow}, together with the \jv{spacedBy} and
\jv{collapsiblePriority} modifier skips --- are gone: those 52 documents are now
in scope, which is where most of the growth from 54 to 99 comes from.

\subsection{Discrepancies found}
\label{sec:discrepancies}

Conformance was not a formality. Six genuine fidelity bugs were found this way,
three while specifying Box, Row and Column and three more while adding the
collapsible and flow managers. In every case the specification was wrong and the
reference implementation was right.

\begin{enumerate}[leftmargin=1.6em]
  \item \textbf{Modifiers modelled as a record.} A flat record of optional
    fields cannot express Section~\ref{sec:orderquirk}. Replacing it with an
    ordered list fixed eight documents.
  \item \textbf{Column size pass modelled as the row's mirror.} The
    \jv{DIRECT\_WEIGHT\_CALCULATION} path of Section~\ref{sec:column} was
    missing.
  \item \textbf{A bare JSON \texttt{weight} read as a horizontal weight.} It is
    \emph{axis-relative}: \jv{DefaultModifierParsers} consults
    \jv{parser.isParentVertical()} and routes \jv{"weight"} to
    \jv{verticalWeight} inside a \jv{column} or \jv{collapsibleColumn} and to
    \jv{horizontalWeight} everywhere else. \jv{collapsiblePriority} resolves its
    orientation by the same rule. Reading the key unconditionally as horizontal
    broke four column documents and twelve checks.
  \item \textbf{A \texttt{spacedBy} \emph{component key} treated as spacing.} No
    component parser reads one; only the modifier of that name reaches a
    container, so a document that declares \jv{"spacedBy"} as a property of the
    component itself gets no spacing at all. Two documents record exactly that.
  \item \textbf{Weight resolution counting collapsed children.} \jv{RowLayout}
    guards both \jv{hasWeights} and \jv{totalWeights} with \jv{isGone}
    (lines 379 and 399), so a child the collapsing sweep dropped neither claims a
    share of the leftover space nor dilutes anyone else's. Modelling the totals
    over all children instead of the visible ones changes every sibling's width.
  \item \textbf{\texttt{applyVisibility} believed to be overridden.} An earlier
    draft of this document asserted that the collapsible managers override it to
    re-run the weight loop. They do not, and nothing else does either; see the
    remark in Section~\ref{sec:row}. The behaviour is unaffected, but the
    justification was wrong.
\end{enumerate}

\begin{remark}[Weighted children of a column can vanish]
Point~3 has a visible consequence that is easy to mistake for a corpus bug. A
child of a \jv{row} authored with \jv{"weight"} takes a share of the row's
leftover width; the same child inside a \jv{column} takes a share of the
column's leftover \emph{height} and keeps a $\Wrap$ width. Over no content that
width is $0$, so the child measures $0 \times h$ and disappears. The gold files
record this, and it is correct: it is what the reference parser does.
\end{remark}

\subsection{Trust boundary}

Not every check has the same status.

\begin{itemize}[leftmargin=1.6em]
  \item The \textbf{properties} of Section~\ref{sec:properties} are proved and
    kernel-checked.
  \item The \textbf{conformance assertions} are decided by kernel reduction of
    \jv{\#guard}. They are as trustworthy as the translation performed by the
    generator.
  \item The \textbf{worked examples} of Appendix~\ref{app:examples} are
    discharged by \jv{native\_decide}, which trusts the compiler as well as the
    kernel.
  \item The correspondence between this specification and the Java is
    established by \emph{reading}, and by the conformance corpus. It is not
    proved, and cannot be.
\end{itemize}

%%% ==================================================================
\appendix

\section{Notation correspondence}

\begin{center}
\small
\begin{longtable}{@{}
  >{\raggedright\arraybackslash}p{0.16\textwidth}
  >{\raggedright\arraybackslash}p{0.36\textwidth}
  >{\raggedright\arraybackslash}p{0.40\textwidth}@{}}
\toprule
\textbf{This document} & \textbf{Lean} & \textbf{Java} \\
\midrule
\endhead
$\Sc$ & \lo{S} & \jv{float} \\
$\Om$ & \lo{fltMax} & \jv{Float.MAX\_VALUE} \\
$p$, $p_{\mathrm{h}}$, $p_{\mathrm{v}}$ & \lo{Padding}, \lo{.horizontal}, \lo{.vertical}
  & \jv{mPaddingLeft} \dots \\
$\kappa$ & \lo{Constraints} & \jv{minWidth}, \jv{maxWidth}, \dots \\
$\kappa_{\mathrm{tight}}$, $\kappa_{\mathrm{loose}}$ & \lo{Constraints.tight}, \lo{.loose} & --- \\
$d$ & \lo{Dim} & \jv{DimensionModifierOperation.Type} \\
$\val$, $\wt$ & \lo{Dim.value}, \lo{Dim.weightValue} & \jv{getValue()}, \jv{hasWeight()} \\
$\base_b$ & \lo{Dim.base} & inline in \jv{computeModifierDefinedWidth} \\
$\iota$ & \lo{DimIn} & \jv{WidthInModifierOperation} \\
$\clampL$, $\clampH$ & \lo{clampLo}, \lo{clampHi} & \jv{Math.max}, \jv{Math.min} \\
$\mu$ & \lo{Modifiers} & \jv{ComponentModifiers} \newline \jv{.getModifiersList()} \\
$\Pp$ & \lo{Modifiers.padding} & \jv{updatePadding} \\
$\Dimw$, $\Dimh$ & \lo{Modifiers.widthDim}, \lo{.heightDim} & \jv{mWidthModifier} \\
$\Iw$, $\Ih$ & \lo{Modifiers.widthIn}, \lo{.heightIn} & \jv{getWidthIn()} \\
$\Dw$, $\Dh$ & \lo{Modifiers.definedWidth}, \lo{.definedHeight} & \jv{computeModifierDefinedWidth} \\
$n$ & \lo{Node} & \jv{Component} subclasses \\
$\tau$ & \lo{LayoutTree} & \jv{ComponentMeasure} in a \jv{MeasurePass} \\
$F_{\mathrm{h}}$, $F_{\mathrm{v}}$ & \lo{Node.isInHorizontalFill}, \lo{.isInVerticalFill}
  & \jv{isInHorizontalFill()} \\
$\Theta_{\mathrm{h}}$, $\Theta_{\mathrm{v}}$ & \lo{totalHWeights}, \lo{totalVWeights} & \jv{totalWeights} \\
$\Rr$ & \lo{resolveAxis} & policy 90--154 \\
$\Mm$ & \lo{measure} & \jv{BaseModernMeasurePolicy.measure} \\
$\mathrm{layout}$ & \lo{layout} & \jv{RootLayoutComponent} \\
$\mathrm{abs}$ & \lo{LayoutTree.absolute} & \jv{getLocationInWindow} \\
$\Aa$ & \lo{arrange} & \jv{internalLayoutMeasure} 465--512 \\
$X$, $Y$ & \lo{crossX}, \lo{crossY} & \jv{internalLayoutMeasure} 524--545 \\
$\mathrm{wrap}_\bullet$ & \lo{boxWrapPass}, \lo{rowWrapPass}, \lo{colWrapPass} & \jv{computeWrapSize} \\
$\mathrm{size}_\bullet$ & \lo{boxSizePass}, \lo{rowSizePass}, \lo{colSizePass} & \jv{computeSize} \\
$\mathrm{place}_\bullet$ & \lo{boxPlace}, \lo{rowPlaceAll}, \lo{colPlaceAll} & \jv{internalLayoutMeasure} \\
\midrule
$g$, $\gone$ & \lo{LayoutTree.gone} & \jv{Component.isGone()} \\
$\miw$, $\mih$ & \lo{minIntrinsicWidth}, \lo{minIntrinsicHeight} & \jv{minIntrinsicWidth} \\
$P_o$ & \lo{Modifiers.priority} & \jv{CollapsiblePriority} \newline \jv{.getPriority} \\
$\prec_o$ & \lo{priorityLt} & the comparator \jv{(int)(p2 - p1)} \\
$\mathrm{sort}_o$ & \lo{sortWithPriorities} & \jv{sortWithPriorities} \\
$\ell_o$ & \lo{findLastStanding} & \jv{findLastStanding} \\
$\Sigma$ & \lo{collapseSweep} & \jv{computeVisibleChildren} 290--325 \\
$G_o$ & \lo{collapseFlags}, \lo{applyFlags} & \jv{setVisibility} \newline \jv{(OVERRIDE\_GONE)} \\
$\Phi$ & \lo{cRowPhase1}, \lo{cColPhase1} & \jv{computeVisibleChildren} 251--281 \\
$e_i$ & \lo{flowDims} & \jv{segmentComponents} 231--253 \\
$\mathrm{seg}$ & \lo{flowSegment} & \jv{segmentComponents} 225--277 \\
$R$, $|r|$ & \lo{flowRowList}, \lo{flowRowMembers} & the \jv{rows} array \\
$\mathrm{mask}_r$ & \lo{maskRow}, \lo{mergeRow} & the per-row \jv{ArrayList} \\
$\eta(r)$ & \lo{flowRowIntrinsic} & \jv{minIntrinsicHeight(ctx, row, true)} \\
\bottomrule
\end{longtable}
\end{center}

\section{Worked examples}
\label{app:examples}

Each example below is stated in the Lean sources as an \jv{example} discharged
by \jv{native\_decide}, so the numbers are computed, not asserted. Rectangles
are listed in pre-order as $(x, y, w, h)$ in window coordinates.

\subsection*{Example 1 --- centring}

A $100 \times 40$ box centred in a $300 \times 200$ viewport:
\[
  \begin{gathered}
    \bbox\bigl(\wdt(\Fill(\bo))\cdot\hgt(\Fill(\bo)),\; \Center,\; \Center,\; [\,c\,]\bigr), \\
    c = \bbox\bigl(\wdt(\Exact(100))\cdot\hgt(\Exact(40)),\; \Start,\; \Top,\; [\,]\bigr)
  \end{gathered}
\]
lays out to $(0,0,300,200)$ and $(100, 80, 100, 40)$. The offsets are
$(300-100)/2 = 100$ and $(200-40)/2 = 80$.

\subsection*{Example 2 --- weights}

A row of three children in a $500 \times 100$ viewport: a fixed $100$ px child
followed by two children of weight $1$ and $3$:
\[
  (0,0,500,100),\quad (0,0,100,20),\quad (100,0,100,20),\quad (200,0,300,20).
\]
The leftover $400$ px is split $1:3$ into $100$ and $300$.

\subsection*{Example 3 --- wrapping with padding}

A wrapping row with $10$ px padding, $s = 5$, and children $30\times30$ and
$40\times20$, nested inside a filling box in a $500\times500$ viewport:
\[
  (0,0,500,500),\quad (0,0,95,50),\quad (0,0,30,30),\quad (35,0,40,20).
\]
The row's content is $30 + 5 + 40 = 75$ wide and $30$ tall, and the padding is
added around it: $10 + 75 + 10 = 95$ and $10 + 30 + 10 = 50$.

The nesting is necessary. A root component is measured against \emph{tight}
viewport constraints, so Step~6 forces even a $\Wrap$ root to the viewport size:
\[
  \mathrm{layout}\bigl(\rrow(\pad(10), \Start, \Top, 0, [\,]),\, 500,\, 500\bigr)
    \text{ has rectangle } (0,0,500,500).
\]

\subsection*{Example 4 --- \texttt{spacedBy} overflows a \texttt{SPACE\_*} arrangement}

A $\SBetween$ row of width $200$, height $50$, with $s = 10$ and two
$50 \times 10$ children:
\[
  (0,0,200,50),\quad (0,0,50,10),\quad (160,0,50,10).
\]
The gap is computed from $T = 100$ as $(200 - 100)/(2-1) = 100$, and then
$s = 10$ is added on top, so the second child starts at $0 + 50 + 100 + 10 =
160$ and ends at $210$ --- ten pixels outside its container. This is the quirk
of Caveat~\ref{sec:spacedquirk}, reproduced exactly.

\subsection*{Example 5 --- a collapsible row drops its lowest priority}

A collapsible row filling a $200 \times 100$ viewport, centred on both axes,
with three $90 \times 50$ children of priorities $1$, $2$ and $3$:
\[
  (0,0,200,100),\quad (10,25,90,50),\quad (10,25,90,50),\quad (100,25,90,50).
\]
\[
  \gone = [\,\textsf{false},\; \textsf{true},\; \textsf{false},\; \textsf{false}\,].
\]
Three children want $270$ px and only $200$ are available, so the sweep, walking
priorities $3, 2, 1$, admits the first two ($90 + 90 \leq 200$) and refuses the
third --- which is the child of priority~$1$, the \emph{first} in declaration
order. The two survivors total $180$, centred at $x = 10$ and $x = 100$.

The collapsed child is nonetheless assigned $x = 10$: it is placed at the cursor
position its successor then also takes, because \jv{setX} runs before the
\jv{isGone} test and only the cursor advance is skipped
(Caveat~\ref{sec:goneplaced}).

\subsection*{Example 6 --- the collapsing fit test ignores \texttt{spacedBy}}

The same shape with two $90 \times 50$ children, $s = 20$, in $180 \times 100$:
\[
  (0,0,180,100),\quad (0,0,90,50),\quad (110,0,90,50),
  \qquad \gone = [\,\textsf{false},\;\textsf{false},\;\textsf{false}\,].
\]
Both children are admitted, because the fit test asks only whether
$90 + 90 \leq 180$. The placement loop then adds the $20$ px gap anyway, so the
second child runs from $110$ to $200$ --- $20$ px past the container it was
just certified to fit inside (Caveat~\ref{sec:collapsespaced}).

\subsection*{Example 7 --- a flow breaks on \texttt{maxColumns}}

A flow with $k_{\mathrm{items}} = 2$ and four $80 \times 40$ children, filling a
$400 \times 300$ viewport:
\[
  (0,0,400,300),\quad (0,0,80,40),\quad (80,0,80,40),\quad
  (0,40,80,40),\quad (80,40,80,40).
\]
All four would fit on one line --- $320 \leq 400$ --- but the item cap cuts a row
after the second. Each row's intrinsic height and its measured height agree here,
both $40$, so the rows stack flush.

\subsection*{Example 8 --- \texttt{maxLines} refuses a child, and it comes back}

A flow with $k_{\mathrm{lines}} = 2$ and five $100 \times 40$ children in a
$260 \times 400$ viewport:
\[
  \begin{gathered}
    (0,0,260,400),\quad (0,0,100,40),\quad (100,0,100,40), \\
    (0,40,100,40),\quad (100,40,100,40),\quad (200,40,100,40),
  \end{gathered}
\]
\[
  \gone = [\,\textsf{false},\;\textsf{false},\;\textsf{false},\;
             \textsf{false},\;\textsf{false},\;\textsf{true}\,].
\]
Two rows of two exhaust the line budget, so the fifth child is refused. That
refusal is written as a \emph{base} visibility, which nothing clears
(Caveat~\ref{sec:flowsticky}); the next pass therefore reads the child back as
zero-width, finds that a zero-width item fits, and lets it rejoin the second row
at the cursor position $x = 200$ --- outside the $260$ px container, and past
both children that genuinely occupy that row. The child stays gone, so nothing
is painted there, but the recorded rectangle is real. This is what
\file{flow\_max\_lines.gold.json} is tagged \jv{suspicious} for, and the
specification reproduces it rather than tidying it away.

\subsection*{Example 9 --- autosize text inside a collapsible column collapses to zero}

A collapsible column of width $300$ and wrapping height, containing a single autosizing
text child with nominal size $42 \times 16.8$, min font size $4$, max font size $100$,
and current font size $14$ (which easily fits within the available height $100$):
\[
  \ccolp\bigl(\wdt(\Exact(300))\cdot\hgt(\Wrap),\; \Center,\; \Center,\; 0,\;
    [\,\ttext(\wdt(\Wrap)\cdot\hgt(\Wrap), \theta_{\mathrm{auto}})\,]\bigr).
\]
Because Phase~1 measures unweighted children with $H^+ = \Om$, the autosizing text balloons
to its maximum font size $100$ with height $120 > 100$. Phase~2 rejects it, yielding:
\[
  \gone = [\,\textsf{true},\; \textsf{true}\,].
\]
Inside a wrapping parent container, the column collapses to height $0$.

\section{Source map}

\begin{center}
\small
\begin{tabular}{@{}llr@{}}
\toprule
\textbf{File} & \textbf{Contents} & \textbf{Lines} \\
\midrule
\file{Basic.lean} & Sections 2--4: scalars, geometry, modifiers & 382 \\
\file{Node.lean} & Sections 3.3, 6.1, 8.4: trees, weights, priorities & 415 \\
\file{Measure.lean} & Sections 5--9: the algorithm & 1134 \\
\file{Theorems.lean} & Section 10: structural properties & 509 \\
\file{Optimizations.lean} & Section 11: optimization passes and invariants & 362 \\
\file{ExtendedModifiers.lean} & Section 12: visibility, baseline, offset, scroll, graphics & 201 \\
\file{ExtendedContainers.lean} & Section 13: FitBox, StateLayout, ImageLayout & 180 \\
\file{WindowAndHitTest.lean} & Section 14: ancestor coordinates, hit testing & 145 \\
\file{Animation.lean} & Section 15: progress, bounds lerp, transitions & 138 \\
\file{DynamicExpressions.lean} & Section 16: expression AST, evaluation, invariants & 145 \\
\file{AutosizeRepaintLoop.lean} & Section 17: Autosize text, false collapse, repaint loop & 266 \\
\file{Conform.lean} & Section 18: traversal and tolerance & 99 \\
\file{Conformance.lean} & Generated: 99 documents, 294 checks & 824 \\
\file{RemoteComposeLayout.lean} & Root module aggregation & 14 \\
\file{tools/gen\_conformance.py} & The corpus translator & 486 \\
\bottomrule
\end{tabular}
\end{center}

\noindent
The Lean modules sit under \file{RemoteComposeLayout/}. Build the development, and with it every check in this document, with
\begin{center}
\file{lake build}
\end{center}
from the \file{lean} directory. Regenerate \file{Conformance.lean} after a
corpus change with \file{python3 tools/gen\_conformance.py}.

\end{document}
